\chapter{Descent over Input Interfaces and Temporal--Interface Interchange}
\label{ch:descent-input-interfaces}

\section{Overview and standing assumptions}
\label{sec:interface-descent-overview}

Chapter~\ref{ch:temporal-trajectories-segal-execution} developed locality in
the temporal variable. For a fixed input interface \(\mathbb A\), a finite
execution restricts to contiguous temporal subintervals and is reconstructed
by gluing compatible adjacent execution segments.

The present chapter studies a second and categorically distinct locality
principle:
\[
    \text{locality in the input interface}.
\]

An internal category
\[
    \mathbb A
\]
describes a structured input environment. Its objects represent input
contexts, modes, configurations, or typing conditions. Its arrows represent
allowed input actions between those contexts, and composition represents
sequential execution of inputs.

An internal functor
\[
    F:\mathbb A'\to\mathbb A
\]
is interpreted as a local presentation of that input environment. Depending on
the application, it may represent:
\begin{enumerate}
    \item one operating mode of a controller;
    \item the commands visible to one subsystem;
    \item one sensor's input vocabulary;
    \item a refinement of global input contexts by local representatives;
    \item a local presentation of a protocol;
    \item a restricted class of disturbances or environmental interactions.
\end{enumerate}

A system
\[
    \Phi:\mathbb A\rightsquigarrow\mathbb B
\]
restricts contravariantly to
\[
    F^\ast\Phi:
    \mathbb A'\rightsquigarrow\mathbb B.
\]
This restriction describes how the same stateful mechanism behaves when it is
exercised only through the local input presentation \(\mathbb A'\).

The corresponding local-to-global question is:

\begin{quote}
Suppose compatible state-transition-output systems have been specified on
overlapping local input presentations. Do they arise from one global system,
and is that global system uniquely determined?
\end{quote}

This is the effective-descent problem for strict Mealy systems.

The overlap compatibility is stronger than equality of observable outputs.
Local systems must identify the same underlying stateful mechanism:
\begin{enumerate}
    \item their state carriers must correspond on overlaps;
    \item corresponding input actions must induce corresponding transitions;
    \item their output comparisons must agree;
    \item the identifications must be coherent on triple overlaps.
\end{enumerate}

Thus interface descent is a theory of patching state machines across
overlapping input vocabularies or input presentations.

Several distinct levels of interface coverage occur:
\begin{enumerate}
    \item coverage of input objects controls reconstruction of state carriers;
    \item coverage of input arrows controls reconstruction of one-step
    transition and output structure;
    \item coverage of composable pairs controls the Mealy multiplication laws;
    \item coverage of all finite input paths controls direct reconstruction of
    finite trajectories;
    \item invertibility of residual restriction controls whether local future
    behaviours are exactly restrictions of global future behaviours;
    \item coverage of all arrows entering locally represented targets controls
    derivative generation and universal safety.
\end{enumerate}

These conditions have different operational meanings and are not
interchangeable. The chapter therefore uses a hierarchy of interface-cover
classes rather than one undifferentiated notion of cover.

The principal results are:
\begin{enumerate}
    \item strict Mealy systems form an input-indexed category;
    \item strict systems satisfy extensive descent over coproduct
    decompositions;
    \item strict systems satisfy effective descent along covers which are
    regular epic in nerve degrees \(0\), \(1\), and \(2\);
    \item canonical residual behaviours restrict strictly along input
    functors;
    \item residual contracts satisfy effective descent along residually exact
    covers;
    \item derivative generation, safety, abstraction, realization, and
    minimization commute with sufficiently exact input restriction;
    \item contract satisfaction, separatedness, and residual reachability are
    interface-local;
    \item finite execution objects satisfy degreewise descent along nervewise
    covers;
    \item temporal gluing and interface reconstruction commute.
\end{enumerate}

Throughout the chapter:
\begin{enumerate}
    \item \(\mathcal E\) is a Grothendieck topos with chosen pullbacks;
    \item \(\mathbb B\) is a fixed internal output category;
    \item \(\mathsf{Inp}\) is a fixed essentially small full subcategory of
    \(\mathsf{IntCat}(\mathcal E)\);
    \item \(\mathsf{Inp}\) is closed under the pullbacks, finite Čech objects,
    and small coproducts used below.
\end{enumerate}

All internal categories, internal functors, Mealy morphisms, strict semantic
homomorphisms, canonical behaviour machines, and trajectory constructions use
the conventions of the preceding chapters.

For an internal functor
\[
    F:\mathbb A'\to\mathbb A,
\]
write
\[
    \overline F
    :=
    F_0\times 1_{B_0}:
    A'_0\times B_0
    \to
    A_0\times B_0.
\]
The notation
\[
    \overline F^{\,\ast}
\]
denotes pullback in the corresponding slices of \(\mathcal E\). The notation
\[
    F^\ast\Phi
\]
denotes input restriction of a Mealy system. Residual restriction and semantic
direct reindexing are denoted respectively by
\[
    F^\sharp
    \qquad\text{and}\qquad
    F^\bullet.
\]

For a covering family
\[
    \{p_i:\mathbb U_i\to\mathbb A\}_{i\in I},
\]
write
\[
    p:=\coprod_i p_i:
    \coprod_i\mathbb U_i\to\mathbb A.
\]
Descent for the family means descent along this coproduct functor. By
extensivity, a system over \(\coprod_i\mathbb U_i\) is equivalently a family
of systems over the \(\mathbb U_i\), and its Čech powers recover the usual
pairwise and triple overlaps.

\section{The input-indexed category of strict systems}
\label{sec:input-indexed-strict-systems}

\subsection{Input restriction}
\label{subsec:input-pullback-strict-systems}

Let
\[
    F:\mathbb A'\to\mathbb A
\]
be an internal functor, and let
\[
    \Phi=(P,\delta_P,\omega_P):
    \mathbb A\rightsquigarrow\mathbb B
\]
be a Mealy morphism.

The carrier of the input restriction is
\[
    P_F
    :=
    A'_0\times_{F_0,A_0,p_P}P.
\]
It fits into the pullback square
\[
\begin{tikzcd}[column sep=large,row sep=large]
    P_F
        \arrow[r, "\pi_P"]
        \arrow[d, "p_{P_F}"']
    &
    P
        \arrow[d, "p_P"]
    \\
    A'_0
        \arrow[r, "F_0"']
    &
    A_0.
\end{tikzcd}
\]
Its output anchor is
\[
    q_{P_F}:=q_P\pi_P.
\]

A generalized element of \(P_F\) is written
\[
    (v',x),
    \qquad
    F_0(v')=p_P(x).
\]

For an arrow
\[
    a':u'\to v'
\]
of \(\mathbb A'\), define
\[
    \delta_{P_F}(a',(v',x))
    :=
    \left(
        u',
        \delta_P(F_1a',x)
    \right),
\]
and
\[
    \omega_{P_F}(a',(v',x))
    :=
    \omega_P(F_1a',x).
\]

\begin{proposition}[Input restriction of strict Mealy systems]
\label{prop:input-pullback-strict-mealy}
The displayed formulas define a Mealy morphism
\[
    F^\ast\Phi:
    \mathbb A'\rightsquigarrow\mathbb B.
\]
If
\[
    f:\Phi\to\Psi
\]
is a strict semantic homomorphism, then
\[
    F^\ast f:
    F^\ast\Phi\to F^\ast\Psi
\]
has carrier map
\[
    (v',x)\longmapsto(v',f(x)).
\]
Consequently, input restriction defines a functor
\[
    F^\ast:
    \mathsf{Mly}^{\mathrm{str}}_{\mathbb A,\mathbb B}
    \to
    \mathsf{Mly}^{\mathrm{str}}_{\mathbb A',\mathbb B}.
\]
\end{proposition}

\begin{proof}
For the transition typing equation,
\[
    p_{P_F}
    \delta_{P_F}(a',(v',x))
    =
    u'
    =
    s_{A'}(a').
\]
For the output typing equations,
\[
    s_B\omega_{P_F}(a',(v',x))
    =
    s_B\omega_P(F_1a',x)
    =
    q_P\delta_P(F_1a',x),
\]
which equals
\[
    q_{P_F}\delta_{P_F}(a',(v',x)).
\]
Similarly,
\[
    t_B\omega_{P_F}(a',(v',x))
    =
    q_{P_F}(v',x).
\]

The unit equations follow from
\[
    F_1(e_{A'}v')=e_A(F_0v')
\]
and the unit equations of \(\Phi\).

For composable arrows
\[
    a':u'\to v',
    \qquad
    b':v'\to w',
\]
functoriality of \(F\) gives
\[
    F_1(b'\circ a')
    =
    F_1b'\circ F_1a'.
\]
The transition and output multiplication equations therefore follow from the
corresponding equations of \(\Phi\).

For a strict homomorphism \(f\), strict transition preservation gives
\[
    f\delta_P(F_1a',x)
    =
    \delta_Q(F_1a',f(x)),
\]
and strict output preservation gives
\[
    \omega_Q(F_1a',f(x))
    =
    \omega_P(F_1a',x).
\]
Thus the pulled-back carrier map is strict.
\end{proof}

\begin{remark}[Systems interpretation of input restriction]
\label{rem:systems-interpretation-input-restriction}
Input restriction does not construct a new physical mechanism. It presents the
same mechanism through a smaller or more refined input vocabulary.

The local state \((v',x)\) records:
\begin{enumerate}
    \item a global machine state \(x\);
    \item a chosen local representative \(v'\) of the global input context
    \(p_P(x)\).
\end{enumerate}
A local input \(a'\) is translated to the global input \(F_1a'\), processed by
the original machine, and then re-anchored at the local source object.

Thus \(F^\ast\Phi\) answers the question:
\[
    \text{how does \(\Phi\) behave when operated through the interface \(F\)?}
\]
\end{remark}

\subsection{Pseudofunctoriality}
\label{subsec:input-pullback-pseudofunctoriality}

For every internal category \(\mathbb A\), there is a canonical isomorphism
\[
    \upsilon_{\mathbb A}:
    1_{\mathbb A}^\ast
    \xrightarrow{\cong}
    1_{\mathsf{Mly}^{\mathrm{str}}_{\mathbb A,\mathbb B}}.
\]

For composable internal functors
\[
\begin{tikzcd}[column sep=large]
    \mathbb A''
        \arrow[r, "G"]
    &
    \mathbb A'
        \arrow[r, "F"]
    &
    \mathbb A,
\end{tikzcd}
\]
associativity of pullbacks gives a canonical isomorphism
\[
    \alpha_{F,G}:
    G^\ast F^\ast
    \xrightarrow{\cong}
    (F\circ G)^\ast.
\]

\begin{proposition}[Coherence of input restriction]
\label{prop:coherence-input-pullback}
The comparison isomorphisms
\[
    \upsilon_{\mathbb A}
    \qquad\text{and}\qquad
    \alpha_{F,G}
\]
satisfy the unit and associativity coherence equations for a pseudofunctor.
\end{proposition}

\begin{proof}
All comparison maps are the canonical maps between iterated pullbacks of the
same carrier diagram. The coherence equations follow from uniqueness in the
universal property of pullbacks.
\end{proof}

\begin{definition}[Input-indexed category of strict systems]
\label{def:input-indexed-system-pseudofunctor}
Define
\[
    \mathsf{Sys}_{\mathbb B}:
    \mathsf{Inp}^{\mathrm{op}}
    \to
    \mathbf{Cat}
\]
by
\[
    \mathsf{Sys}_{\mathbb B}(\mathbb A)
    :=
    \mathsf{Mly}^{\mathrm{str}}_{\mathbb A,\mathbb B},
\]
and
\[
    \mathsf{Sys}_{\mathbb B}(F)
    :=
    F^\ast.
\]
The coherence isomorphisms are those of
Proposition~\ref{prop:coherence-input-pullback}.
\end{definition}

\subsection{The total category of systems}
\label{subsec:grothendieck-construction-input-systems}

\begin{definition}[Total category of strict systems]
\label{def:total-category-strict-systems}
The Grothendieck construction
\[
    \int\mathsf{Sys}_{\mathbb B}
\]
has:
\begin{enumerate}
    \item objects \((\mathbb A,\Phi)\), where
    \[
        \mathbb A\in\mathsf{Inp}
    \]
    and
    \[
        \Phi:\mathbb A\rightsquigarrow\mathbb B;
    \]
    \item morphisms
    \[
        (\mathbb A',\Phi')
        \longrightarrow
        (\mathbb A,\Phi)
    \]
    consisting of an internal functor
    \[
        F:\mathbb A'\to\mathbb A
    \]
    and a strict semantic homomorphism
    \[
        \varphi:
        \Phi'\to F^\ast\Phi.
    \]
\end{enumerate}
\end{definition}

Composition uses the comparison
\[
\begin{tikzcd}[column sep=large]
    G^\ast F^\ast\Phi
        \arrow[r, "\alpha_{F,G}", "\cong"']
    &
    (F\circ G)^\ast\Phi.
\end{tikzcd}
\]

The projection
\[
    \pi_{\mathsf{Sys}}:
    \int\mathsf{Sys}_{\mathbb B}
    \to
    \mathsf{Inp}
\]
sends
\[
    (\mathbb A,\Phi)\longmapsto\mathbb A
\]
and
\[
    (F,\varphi)\longmapsto F.
\]

\begin{proposition}[Fibration of strict systems]
\label{prop:fibration-strict-systems}
The projection
\[
    \pi_{\mathsf{Sys}}
\]
is a Grothendieck fibration. A morphism over
\[
    F:\mathbb A'\to\mathbb A
\]
is cartesian precisely when its strict component is an isomorphism
\[
    \Phi'\xrightarrow{\cong}F^\ast\Phi.
\]
\end{proposition}

\begin{proof}
For every \((\mathbb A,\Phi)\) and every
\(F:\mathbb A'\to\mathbb A\), the morphism
\[
    (\mathbb A',F^\ast\Phi)
    \longrightarrow
    (\mathbb A,\Phi)
\]
with strict component
\[
    1_{F^\ast\Phi}
\]
has the cartesian universal property. The characterization follows from the
standard Grothendieck construction of a contravariant pseudofunctor.
\end{proof}

\section{Čech descent for strict systems}
\label{sec:cech-descent-strict-systems}

\subsection{Čech powers of an input functor}
\label{subsec:cech-powers-input-functor}

Let
\[
    p:\mathbb U\to\mathbb A
\]
be an internal functor.

Define the Čech powers
\[
    \mathbb U^{\langle r\rangle}
    :=
    \underbrace{
    \mathbb U\times_{\mathbb A}
    \cdots\times_{\mathbb A}\mathbb U
    }_{r\text{ factors}}
\]
for \(r\geq1\).

In particular,
\[
\begin{tikzcd}[column sep=large]
    \mathbb U^{\langle2\rangle}
        =
    \mathbb U\times_{\mathbb A}\mathbb U
        \arrow[r, shift left=0.7ex, "\pi_1"]
        \arrow[r, shift right=0.7ex, "\pi_2"']
    &
    \mathbb U,
\end{tikzcd}
\]
and
\[
    \mathbb U^{\langle3\rangle}
    =
    \mathbb U\times_{\mathbb A}
    \mathbb U\times_{\mathbb A}\mathbb U.
\]

Write
\[
    \pi_{12},\pi_{23},\pi_{13}:
    \mathbb U^{\langle3\rangle}
    \to
    \mathbb U^{\langle2\rangle}
\]
for the three pairwise projections.

The overlap category
\[
    \mathbb U^{\langle2\rangle}
\]
records two local presentations with the same global image. The triple
overlap records three such presentations simultaneously.

\subsection{System descent structures}
\label{subsec:system-descent-data-rigorous}

\begin{definition}[System descent structure]
\label{def:system-descent-datum-rigorous}
A \textbf{system descent structure} for
\[
    p:\mathbb U\to\mathbb A
\]
consists of:
\begin{enumerate}
    \item a strict Mealy system
    \[
        \Phi_{\mathbb U}
        \in
        \mathsf{Sys}_{\mathbb B}(\mathbb U);
    \]
    \item an isomorphism
    \[
        \vartheta:
        \pi_1^\ast\Phi_{\mathbb U}
        \xrightarrow{\cong}
        \pi_2^\ast\Phi_{\mathbb U}
    \]
    over
    \(\mathbb U^{\langle2\rangle}\);
    \item the cocycle equation
    \[
        \pi_{23}^\ast\vartheta
        \circ
        \pi_{12}^\ast\vartheta
        =
        \pi_{13}^\ast\vartheta
    \]
    over
    \(\mathbb U^{\langle3\rangle}\),
    after inserting the canonical pseudofunctorial comparison
    isomorphisms.
\end{enumerate}
\end{definition}

A morphism
\[
    f:
    (\Phi_{\mathbb U},\vartheta)
    \to
    (\Psi_{\mathbb U},\chi)
\]
of system descent structures is a strict semantic homomorphism
\[
    f:\Phi_{\mathbb U}\to\Psi_{\mathbb U}
\]
such that
\[
    \chi\circ\pi_1^\ast f
    =
    \pi_2^\ast f\circ\vartheta.
\]

Write
\[
    \mathsf{Desc}_p
    \left(
        \mathsf{Sys}_{\mathbb B}
    \right)
\]
for the category of system descent structures.

\begin{remark}[Operational meaning of overlap coherence]
\label{rem:operational-overlap-coherence}
The isomorphism \(\vartheta\) says that two local presentations of the same
global input situation describe the same stateful mechanism.

It identifies:
\begin{enumerate}
    \item the corresponding local states;
    \item the transitions produced by corresponding input actions;
    \item the associated output comparisons.
\end{enumerate}

The cocycle equation says that changing from local presentation \(1\) to
presentation \(2\), and then from \(2\) to \(3\), gives the same state
identification as changing directly from \(1\) to \(3\).

Thus pairwise compatibility is not sufficient. The local coordinate changes
for the state machine must be jointly coherent.
\end{remark}

\subsection{The comparison functor}
\label{subsec:system-descent-comparison}

Every global system
\[
    \Phi\in\mathsf{Sys}_{\mathbb B}(\mathbb A)
\]
determines a canonical system descent structure. Its local system is
\[
    p^\ast\Phi.
\]
The two pullbacks to
\(\mathbb U^{\langle2\rangle}\) are canonically isomorphic because
\[
    p\pi_1=p\pi_2.
\]

\begin{definition}[System descent comparison]
\label{def:system-descent-comparison-rigorous}
The \textbf{system descent comparison functor} is
\[
    K_p:
    \mathsf{Sys}_{\mathbb B}(\mathbb A)
    \to
    \mathsf{Desc}_p
    \left(
        \mathsf{Sys}_{\mathbb B}
    \right).
\]
\end{definition}

\subsection{Identifiability and integrability}
\label{subsec:identifiability-integrability-categorical}

\begin{definition}[Interface identifiability]
\label{def:interface-identifiability-rigorous}
The input functor \(p\) is \textbf{identifying for strict systems} if
\[
    K_p
\]
is fully faithful.
\end{definition}

\begin{definition}[Interface integrability]
\label{def:interface-integrability-rigorous}
The input functor \(p\) is \textbf{integrating for strict systems} if
\[
    K_p
\]
is essentially surjective.
\end{definition}

\begin{definition}[Effective system descent]
\label{def:effective-system-descent}
The input functor \(p\) has \textbf{effective descent for strict systems} if
\[
    K_p
\]
is an equivalence.
\end{definition}

Identifiability means that no globally meaningful distinction between systems
or strict system maps is invisible on every local patch.

Integrability means that every compatible local family can be implemented by
one global system.

Effective descent combines these two statements:
\[
\begin{tikzcd}[column sep=huge]
    \text{compatible local systems}
        \arrow[r, leftrightarrow, "\text{effective descent}"]
    &
    \text{global system, unique up to isomorphism}.
\end{tikzcd}
\]

\section{System cover classes}
\label{sec:system-cover-classes}

\subsection{Nerve components}
\label{subsec:nerve-components-input-cover}

For an internal category \(\mathbb A\), write
\[
    A_n:=N\mathbb A_n.
\]
Thus
\[
    A_0
\]
is the object object,
\[
    A_1
\]
is the arrow object, and
\[
    A_2
    =
    A_1\times_{t_A,A_0,s_A}A_1
\]
is the object of composable pairs.

An internal functor
\[
    p:\mathbb U\to\mathbb A
\]
induces maps
\[
    p_n:
    U_n\to A_n
\]
for every \(n\geq0\).

\subsection{Extensive covers}
\label{subsec:extensive-interface-covers}

\begin{definition}[Extensive interface cover]
\label{def:extensive-interface-cover}
An \textbf{extensive interface cover} is a family of coproduct injections
\[
    \iota_i:
    \mathbb A_i\to
    \coprod_{i\in I}\mathbb A_i.
\]
\end{definition}

The object and arrow objects decompose as
\[
    \left(
        \coprod_i\mathbb A_i
    \right)_0
    =
    \coprod_iA_{i,0},
\]
and
\[
    \left(
        \coprod_i\mathbb A_i
    \right)_1
    =
    \coprod_iA_{i,1}.
\]
There are no arrows between distinct components.

Systems-theoretically, an extensive cover represents a genuinely disconnected
input environment. A system on the whole interface is therefore exactly an
independent family of systems on its connected components.

\subsection{\texorpdfstring{\(2\)}{2}-nerve covers}
\label{subsec:two-nerve-covers}

\begin{definition}[\(2\)-nerve cover]
\label{def:two-nerve-cover}
An internal functor
\[
    p:\mathbb U\to\mathbb A
\]
is a \textbf{\(2\)-nerve cover} if
\[
    p_0:U_0\twoheadrightarrow A_0,
\]
\[
    p_1:U_1\twoheadrightarrow A_1,
\]
and
\[
    p_2:U_2\twoheadrightarrow A_2
\]
are regular epimorphisms.
\end{definition}

The three conditions have distinct operational meanings.

\paragraph{Degree zero.}
Every global input context is locally represented. This is what allows states
anchored at every global context to be reconstructed.

\paragraph{Degree one.}
Every global input action is locally represented. This allows reconstruction
of the one-step transition and output response.

\paragraph{Degree two.}
Every globally executable pair of inputs is jointly represented. This allows
one to check that processing two inputs successively agrees with processing
their composite.

The distinction is important. A collection of local interfaces may contain
every individual command while failing to contain some executable pair in one
common patch. In that situation, the response to each command is visible, but
their sequential interaction is not locally testable.

\subsection{Nervewise covers}
\label{subsec:nervewise-covers}

\begin{definition}[Nervewise interface cover]
\label{def:nervewise-interface-cover}
An internal functor
\[
    p:\mathbb U\to\mathbb A
\]
is a \textbf{nervewise cover} if
\[
    p_n:U_n\twoheadrightarrow A_n
\]
is a regular epimorphism for every \(n\geq0\).
\end{definition}

Every nervewise cover is a \(2\)-nerve cover. The converse need not hold.

A \(2\)-nerve cover reconstructs the rules of the machine. A nervewise cover
additionally represents every finite input history and therefore supports
direct descent of already-expanded finite runs.

\subsection{Stability of nerve covers}
\label{subsec:stability-nerve-covers}

\begin{proposition}[Stability of nerve covers]
\label{prop:stability-nerve-covers}
The classes of \(2\)-nerve covers and nervewise covers:
\begin{enumerate}
    \item contain all identity functors;
    \item are closed under composition;
    \item are stable under pullback in \(\mathsf{Inp}\).
\end{enumerate}
\end{proposition}

\begin{proof}
The nerve functor preserves limits, so a pullback of internal functors induces
a pullback in every nerve degree. Regular epimorphisms in a Grothendieck topos
contain identities and are stable under pullback and composition.
\end{proof}

\subsection{System and trajectory topologies}
\label{subsec:system-trajectory-topologies}

\begin{definition}[System interface pretopology]
\label{def:system-interface-pretopology}
Let
\[
    \mathcal K_{\mathrm{sys}}
\]
be the smallest Grothendieck pretopology on \(\mathsf{Inp}\) containing:
\begin{enumerate}
    \item all extensive interface covers;
    \item all singleton \(2\)-nerve covers.
\end{enumerate}
Write
\[
    \tau_{\mathrm{sys}}
\]
for the generated Grothendieck topology.
\end{definition}

\begin{definition}[Trajectory interface pretopology]
\label{def:trajectory-interface-pretopology}
Let
\[
    \mathcal K_{\mathrm{traj}}
\]
be the smallest Grothendieck pretopology on \(\mathsf{Inp}\) containing:
\begin{enumerate}
    \item all extensive interface covers;
    \item all singleton nervewise covers.
\end{enumerate}
Write
\[
    \tau_{\mathrm{traj}}
\]
for the generated Grothendieck topology.
\end{definition}

\section{Extensive descent}
\label{sec:extensive-descent-input-interfaces}

\subsection{Decomposition of carriers}
\label{subsec:extensive-decomposition-carriers}

Let
\[
    \mathbb A
    =
    \coprod_{i\in I}\mathbb A_i.
\]
Then
\[
    A_0=\coprod_iA_{i,0}.
\]

Since \(\mathcal E\) is infinitary extensive, there is a canonical equivalence
\[
    \mathcal E/
    \left(
        \left(
            \coprod_iA_{i,0}
        \right)
        \times B_0
    \right)
    \simeq
    \prod_i
    \mathcal E/(A_{i,0}\times B_0).
\]

A carrier
\[
    P\to
    \left(
        \coprod_iA_{i,0}
    \right)
    \times B_0
\]
therefore decomposes, up to canonical isomorphism, as
\[
    P\cong\coprod_iP_i
\]
with
\[
    P_i\to A_{i,0}\times B_0.
\]

\subsection{Decomposition of transition domains}
\label{subsec:extensive-decomposition-derivative-domains}

Since
\[
    A_1=\coprod_iA_{i,1},
\]
and every arrow belongs to a unique component,
\[
    A_1\times_{A_0}P
    \cong
    \coprod_i
    \left(
        A_{i,1}\times_{A_{i,0}}P_i
    \right).
\]

Consequently:
\begin{enumerate}
    \item a transition map on \(P\) is equivalent to a family of transition
    maps on the \(P_i\);
    \item an output map on \(P\) is equivalent to a family of output maps on
    the \(P_i\);
    \item every typing, unit, and multiplication equation is equivalent to its
    componentwise family.
\end{enumerate}

\begin{theorem}[Extensive descent for strict Mealy systems]
\label{thm:extensive-descent-strict-systems}
There is a canonical equivalence
\[
\begin{tikzcd}[column sep=huge]
    \mathsf{Sys}_{\mathbb B}
    \left(
        \coprod_{i\in I}\mathbb A_i
    \right)
        \arrow[r, "\simeq"]
    &
    \prod_{i\in I}
    \mathsf{Sys}_{\mathbb B}(\mathbb A_i).
\end{tikzcd}
\]
\end{theorem}

\begin{proof}
Infinitary extensivity gives the equivalence on carriers. The transition
domain decomposes as the coproduct of the component transition domains.
Transition and output maps are therefore exactly componentwise families, and
all Mealy equations are componentwise. Strict semantic homomorphisms decompose
in the same manner.
\end{proof}

\subsection{Decomposition of canonical behaviours}
\label{subsec:extensive-decomposition-canonical-behaviours}

Let
\[
    v\in A_{i,0}.
\]
Because there are no arrows between distinct coproduct components, there is a
canonical isomorphism
\[
    \left(
        \coprod_j\mathbb A_j
    \right)/v
    \cong
    \mathbb A_i/v.
\]

\begin{proposition}[Extensive decomposition of canonical behaviours]
\label{prop:extensive-decomposition-canonical-behaviours}
There is a canonical isomorphism over
\[
    \left(
        \coprod_iA_{i,0}
    \right)
    \times B_0
\]
of canonical behaviour carriers
\[
    \mathbf{Beh}_{\coprod_i\mathbb A_i,\mathbb B,0}
    \cong
    \coprod_i
    \mathbf{Beh}_{\mathbb A_i,\mathbb B,0}.
\]
This isomorphism preserves the canonical transition and output structure.
\end{proposition}

\begin{proof}
Over an object \(v\in A_{i,0}\), a residual behaviour is an internal functor
\[
    \left(
        \coprod_j\mathbb A_j
    \right)/v
    \to
    \mathbb B.
\]
The residual-category isomorphism identifies this with a functor
\[
    \mathbb A_i/v\to\mathbb B.
\]
The identification is natural in residual restriction and preserves
evaluation of output comparisons. It therefore assembles into an isomorphism
of canonical behaviour Mealy morphisms.
\end{proof}

\begin{proposition}[Residual exactness of coproduct inclusions]
\label{prop:residual-exactness-coproduct-inclusions}
For every coproduct inclusion
\[
    \iota_i:
    \mathbb A_i\to\coprod_j\mathbb A_j,
\]
the functor \(\iota_i\) is universally behaviour-exact and universally
derivative-covering.
\end{proposition}

\begin{proof}
For every \(v\in A_{i,0}\), there is a canonical isomorphism
\[
    \left(\coprod_j\mathbb A_j\right)/v
    \cong
    \mathbb A_i/v.
\]
Residual restriction along \(\iota_i\) is therefore an isomorphism.

Every arrow of the coproduct category whose target lies in
\(\mathbb A_i\) belongs uniquely to \(\mathbb A_i\). Hence the
incoming-arrow comparison for \(\iota_i\) is an isomorphism.

To check the universal clauses, let
\[
    H:\mathbb C\to\coprod_j\mathbb A_j
\]
be any internal functor. The object object \(C_0\) decomposes according to
which coproduct component contains its image. Since the source and target of
an arrow have images in the same component, \(\mathbb C\) decomposes as the
coproduct of the corresponding full inverse-image categories. Thus a pullback
of \(\iota_i\) is again a coproduct-component inclusion. The preceding slice
and incoming-arrow arguments apply after every base change.
\end{proof}

\subsection{Extensive descent of contracts}
\label{subsec:extensive-descent-residual-contracts}

\begin{theorem}[Extensive descent for residual contracts]
\label{thm:extensive-descent-residual-contracts}
There is a canonical isomorphism of posets
\[
    \mathsf{Ctr}_{\mathbb B}
    \left(
        \coprod_i\mathbb A_i
    \right)
    \cong
    \prod_i
    \mathsf{Ctr}_{\mathbb B}(\mathbb A_i).
\]
\end{theorem}

\begin{proof}
By
Proposition~\ref{prop:extensive-decomposition-canonical-behaviours},
subobjects of the global canonical behaviour carrier decompose componentwise.
Every input arrow lies in a unique component, so derivative invariance is
equivalent to componentwise derivative invariance.
\end{proof}

\subsection{Extensive descent of temporal trajectories}
\label{subsec:extensive-descent-temporal-trajectories}

Let
\[
    \Phi
    \cong
    \coprod_i\Phi_i
\]
under
Theorem~\ref{thm:extensive-descent-strict-systems}.

\begin{proposition}[Extensive decomposition of execution categories]
\label{prop:extensive-decomposition-execution-categories}
There is a canonical isomorphism
\[
    \operatorname{Exec}(\Phi)
    \cong
    \coprod_i
    \operatorname{Exec}(\Phi_i).
\]
\end{proposition}

\begin{proof}
The object object is
\[
    P\cong\coprod_iP_i.
\]
The arrow object is
\[
    A_1\times_{A_0}P
    \cong
    \coprod_i
    \left(
        A_{i,1}\times_{A_{i,0}}P_i
    \right).
\]
All source, target, identity, and composition maps are componentwise.
\end{proof}

\begin{corollary}[Extensive descent of trajectories]
\label{cor:extensive-descent-trajectories}
There is a canonical isomorphism of trajectory sheaves
\[
    \mathscr T_\Phi
    \cong
    \coprod_i\mathscr T_{\Phi_i}.
\]
\end{corollary}

\begin{proof}
Apply the trajectory functor
\[
    \operatorname{Path}
\]
to
Proposition~\ref{prop:extensive-decomposition-execution-categories}.
Finite paths cannot cross between distinct coproduct components.
\end{proof}

\section{Effective descent of strict Mealy systems}
\label{sec:effective-descent-strict-mealy-systems}

Throughout this section, let
\[
    p:\mathbb U\to\mathbb A
\]
be a \(2\)-nerve cover.

\subsection{Descent of carriers}
\label{subsec:descent-of-carriers}

Let
\[
    \Phi_{\mathbb U}
    =
    (P_{\mathbb U},
    \delta_{\mathbb U},
    \omega_{\mathbb U})
\]
be the local system of a system descent structure.

Its carrier is an object
\[
    P_{\mathbb U}
    \to
    U_0\times B_0.
\]
The carrier component of the overlap isomorphism gives ordinary descent
structure along
\[
    p_0\times1_{B_0}:
    U_0\times B_0
    \twoheadrightarrow
    A_0\times B_0.
\]

Regular epimorphisms in a Grothendieck topos are effective descent morphisms.
Hence there exists an object
\[
    P\to A_0\times B_0
\]
and an isomorphism
\[
    \xi_P:
    \overline p^{\,\ast}P
    \xrightarrow{\cong}
    P_{\mathbb U}
\]
compatible with the carrier overlap identification.

Write
\[
    p_P:P\to A_0
\]
and
\[
    q_P:P\to B_0
\]
for the two anchors.

Operationally, surjectivity of \(p_0\) means that every global operating
context occurs in at least one local presentation. Without this condition,
states anchored over a missing context could not be recovered from the local
systems.

\subsection{Transition-domain base change}
\label{subsec:derivative-domain-base-change}

Define the global transition domain
\[
    D_P
    :=
    A_1\times_{t_A,A_0,p_P}P.
\]
The local transition domain is
\[
    D_{P_{\mathbb U}}
    :=
    U_1
    \times_{t_U,U_0,p_{P_{\mathbb U}}}
    P_{\mathbb U}.
\]

\begin{lemma}[Transition domains commute with input restriction]
\label{lem:derivative-domains-input-pullback}
Under the carrier identification
\[
    P_{\mathbb U}
    \cong
    U_0\times_{p_0,A_0,p_P}P,
\]
there is a canonical isomorphism over \(U_1\)
\[
    D_{P_{\mathbb U}}
    \cong
    U_1\times_{p_1,A_1}D_P.
\]
\end{lemma}

\begin{proof}
The right-hand side consists of triples
\[
    (a',a,x)
\]
such that
\[
    a=p_1(a')
\]
and
\[
    t_A(a)=p_P(x).
\]
Functoriality of \(p\) gives
\[
    t_Ap_1(a')
    =
    p_0t_U(a').
\]
Thus the same structure is equivalently a pair
\[
    (a',(t_Ua',x))
\]
with
\[
    p_0t_U(a')=p_P(x),
\]
which is precisely an element of
\[
    U_1\times_{U_0}P_{\mathbb U}.
\]
The comparison is the canonical pullback-pasting isomorphism.
\end{proof}

\subsection{Compatibility of the local operation maps}
\label{subsec:compatibility-local-operation-maps}

The nerve preserves pullbacks. Hence
\[
    N(\mathbb U\times_{\mathbb A}\mathbb U)_1
    \cong
    U_1\times_{A_1}U_1,
\]
and
\[
    N(\mathbb U\times_{\mathbb A}\mathbb U)_2
    \cong
    U_2\times_{A_2}U_2.
\]

The strictness of the overlap isomorphism
\[
    \vartheta:
    \pi_1^\ast\Phi_{\mathbb U}
    \xrightarrow{\cong}
    \pi_2^\ast\Phi_{\mathbb U}
\]
therefore says exactly that:
\begin{enumerate}
    \item the two pullbacks of the local transition map to
    \(U_1\times_{A_1}U_1\) agree after transporting states through
    \(\vartheta\);
    \item the corresponding output maps agree;
    \item the induced comparisons on composable pairs are coherent.
\end{enumerate}

Thus the local operation maps are morphisms between effective descent objects
over \(A_1\), rather than merely maps whose domains and codomains happen to
descend.

\subsection{Descent of transition maps}
\label{subsec:descent-transition-maps-rigorous}

The transition typing equation
\[
    p_{P_{\mathbb U}}\delta_{\mathbb U}
    =
    s_U\pi_{U_1}
\]
allows the local transition to be regarded as a map in
\(\mathcal E/U_1\):
\[
    \widetilde\delta_{\mathbb U}:
    D_{P_{\mathbb U}}
    \to
    s_U^\ast P_{\mathbb U}.
\]

Under
Lemma~\ref{lem:derivative-domains-input-pullback},
both source and target are pullbacks along
\[
    p_1:U_1\twoheadrightarrow A_1
\]
of
\[
    D_P
    \qquad\text{and}\qquad
    s_A^\ast P.
\]

\begin{proposition}[Descent of transition structure]
\label{prop:descent-transition-structure}
There exists a unique map over \(A_1\)
\[
    \widetilde\delta_P:
    D_P\to s_A^\ast P
\]
whose pullback along \(p_1\) is
\(\widetilde\delta_{\mathbb U}\).
Its composite with
\[
    s_A^\ast P\to P
\]
defines a global transition map
\[
    \delta_P:D_P\to P.
\]
\end{proposition}

\begin{proof}
The map
\[
    p_1:U_1\twoheadrightarrow A_1
\]
is an effective descent morphism. The overlap compatibility of
\(\widetilde\delta_{\mathbb U}\) was identified in the preceding subsection.
Morphisms between descended objects therefore descend uniquely.
\end{proof}

\subsection{Descent of output maps}
\label{subsec:descent-output-maps-rigorous}

Regard the local output map as a map over \(U_1\):
\[
    \widetilde\omega_{\mathbb U}
    :=
    \left\langle
        \pi_{U_1},
        \omega_{\mathbb U}
    \right\rangle:
    D_{P_{\mathbb U}}
    \to
    U_1\times B_1.
\]

The object \(U_1\times B_1\) is the pullback of
\[
    A_1\times B_1
\]
along \(p_1\).

\begin{proposition}[Descent of output structure]
\label{prop:descent-output-structure}
There exists a unique map over \(A_1\)
\[
    \widetilde\omega_P:
    D_P\to A_1\times B_1
\]
whose pullback is
\(\widetilde\omega_{\mathbb U}\).
Projection to \(B_1\) defines
\[
    \omega_P:D_P\to B_1.
\]
\end{proposition}

\begin{proof}
The overlap compatibility follows from strictness of the local system
identification. Effective descent along \(p_1\) therefore descends the map
uniquely.
\end{proof}

\subsection{Reflection of typing equations}
\label{subsec:reflection-typing-equations}

The transition typing equation holds by construction. The remaining output
typing equations are
\[
    s_B\omega_P
    =
    q_P\delta_P,
\]
and
\[
    t_B\omega_P
    =
    q_P\pi_P.
\]

\begin{proposition}[Descent of typing equations]
\label{prop:descent-typing-equations}
The descended transition and output maps satisfy all Mealy typing equations.
\end{proposition}

\begin{proof}
Pulling the candidate equations back along
\[
    p_1:U_1\twoheadrightarrow A_1
\]
gives the local typing equations. Pullback along a regular epimorphism reflects
equality of maps, so the global equations hold.
\end{proof}

\subsection{Reflection of unit equations}
\label{subsec:reflection-unit-equations}

\begin{proposition}[Descent of unit equations]
\label{prop:descent-unit-equations}
The descended transition and output maps satisfy the Mealy unit laws.
\end{proposition}

\begin{proof}
The transition unit equation is an equality of maps
\[
    P\rightrightarrows P,
\]
and the output unit equation is an equality of maps
\[
    P\rightrightarrows B_1.
\]
Pulling them back along
\[
    \overline p:
    U_0\times B_0
    \twoheadrightarrow
    A_0\times B_0
\]
gives the local unit equations. Since \(\overline p\) is a regular
epimorphism, pullback reflects equality.
\end{proof}

\subsection{Reflection of multiplication equations}
\label{subsec:reflection-multiplication-equations}

Let
\[
    A_2
    =
    A_1\times_{t_A,A_0,s_A}A_1
\]
be the object of composable pairs
\[
    a:u\to v,
    \qquad
    b:v\to w.
\]
Write
\[
    \pi_1,\pi_2:
    A_2\rightrightarrows A_1
\]
for the first and second arrows.

Define the two-step execution domain
\[
    D_P^{(2)}
    :=
    A_2
    \times_{t_A\pi_2,A_0,p_P}
    P.
\]

The transition multiplication law is an equality
\[
    D_P^{(2)}
    \rightrightarrows
    P,
\]
and the output multiplication law is an equality
\[
    D_P^{(2)}
    \rightrightarrows
    B_1.
\]

\begin{proposition}[Descent of multiplication equations]
\label{prop:descent-multiplication-equations}
The descended transition and output maps satisfy the Mealy multiplication
laws.
\end{proposition}

\begin{proof}
Pulling
\[
    D_P^{(2)}
\]
back along
\[
    p_2:U_2\twoheadrightarrow A_2
\]
gives the local two-step execution domain. Under this identification, the
pullbacks of the two global transition expressions are the two sides of the
local transition multiplication law. The same holds for the output
multiplication expressions.

Since \(p_2\) is a regular epimorphism, pullback along \(p_2\) reflects
equality.
\end{proof}

\begin{remark}[Why composable pairs must be covered]
\label{rem:why-composable-pairs-covered}
Coverage of individual arrows determines the response to each input
separately. It does not, by itself, determine whether two locally specified
responses fit the global sequential law.

Coverage of \(A_2\) says that every globally executable pair is visible
together in some local presentation. It is therefore the exact level at which
one can verify
\[
    \text{response to a composite}
    =
    \text{successive responses}.
\]
\end{remark}

\subsection{Descent of strict homomorphisms}
\label{subsec:descent-strict-homomorphisms}

Let
\[
    (\Phi_{\mathbb U},\vartheta)
\]
and
\[
    (\Psi_{\mathbb U},\chi)
\]
be system descent structures, and let
\[
    f_{\mathbb U}:
    \Phi_{\mathbb U}\to\Psi_{\mathbb U}
\]
be a compatible local strict homomorphism.

Its carrier map descends uniquely along
\[
    \overline p
\]
to
\[
    f:P\to Q.
\]

\begin{proposition}[Descent of strict semantic homomorphisms]
\label{prop:descent-strict-semantic-homomorphisms}
The descended carrier map \(f\) is a strict semantic homomorphism.
\end{proposition}

\begin{proof}
The transition-preservation and output-preservation equations are equalities
of maps over \(A_1\). Their pullbacks along \(p_1\) are the strict
preservation equations for \(f_{\mathbb U}\). Since \(p_1\) is a regular
epimorphism, the global equations hold.
\end{proof}

\subsection{The effective descent theorem}
\label{subsec:effective-descent-theorem-strict-systems}

\begin{theorem}[Effective descent for strict Mealy systems]
\label{thm:effective-descent-strict-mealy-systems}
For every \(2\)-nerve cover
\[
    p:\mathbb U\to\mathbb A,
\]
the comparison functor
\[
    K_p:
    \mathsf{Sys}_{\mathbb B}(\mathbb A)
    \longrightarrow
    \mathsf{Desc}_p
    \left(
        \mathsf{Sys}_{\mathbb B}
    \right)
\]
is an equivalence.
\end{theorem}

\begin{proof}
\leavevmode
\begin{enumerate}
    \item \textbf{Essential surjectivity.}
    Given a system descent structure, descend the carrier along \(p_0\), the
    transition and output maps along \(p_1\), and reflect the typing, unit, and
    multiplication laws along \(p_1\), \(p_0\), and \(p_2\), respectively.
    The preceding propositions produce a global strict Mealy system whose
    pullback is isomorphic to the local system with its overlap
    identification.

    \item \textbf{Fullness.}
    A compatible local strict homomorphism descends on carriers along \(p_0\).
    Proposition
    \ref{prop:descent-strict-semantic-homomorphisms}
    shows that the descended map is strict.

    \item \textbf{Faithfulness.}
    If two global strict homomorphisms have equal pullbacks, their carrier maps
    become equal after pullback along the regular epimorphism
    \(\overline p\). Hence they are equal globally.
\end{enumerate}
\end{proof}

\begin{remark}[Systems-theoretic form of the theorem]
\label{rem:systems-form-effective-descent}
The theorem says that a global state machine is reconstructed from local
models provided:
\begin{enumerate}
    \item every input context is locally represented;
    \item every one-step input action is locally represented;
    \item every executable pair of input actions is locally represented;
    \item the local state, transition, and output models agree coherently on
    overlaps.
\end{enumerate}

No additional coverage of longer paths is needed to reconstruct the machine
itself, because its longer executions are generated by one-step transition
structure and the multiplication laws.
\end{remark}

\subsection{The stack of strict systems}
\label{subsec:stack-of-strict-systems}

\begin{lemma}[Composition of effective descent]
\label{lem:composition-effective-descent}
Let
\[
    \mathcal F:
    \mathcal C^{\mathrm{op}}\to\mathbf{Cat}
\]
be a pseudofunctor. Suppose \(\mathcal F\) satisfies effective descent for a
first covering family and for the pullbacks of a second covering family to the
domains of the first. Then it satisfies effective descent for the composite
covering family.
\end{lemma}

\begin{proof}
A compatible local family for the composite cover is equivalently:
\begin{enumerate}
    \item a compatible family on each first-stage domain for the pulled-back
    second-stage cover;
    \item compatibility among the resulting first-stage objects.
\end{enumerate}
The first collection of descent equivalences reconstructs the first-stage
objects. The second reconstructs the global object. The same argument applies
to morphisms.
\end{proof}

\begin{corollary}[Stack of strict Mealy systems]
\label{cor:stack-strict-mealy-systems}
The pseudofunctor
\[
    \mathsf{Sys}_{\mathbb B}
\]
is a stack for
\[
    \tau_{\mathrm{sys}}.
\]
\end{corollary}

\begin{proof}
Effective descent for extensive covers is
Theorem~\ref{thm:extensive-descent-strict-systems}.
Effective descent for singleton \(2\)-nerve covers is
Theorem~\ref{thm:effective-descent-strict-mealy-systems}.
The generating cover classes are stable under pullback, and effective descent
is stable under composition by
Lemma~\ref{lem:composition-effective-descent}.
\end{proof}

\section{Residual restriction and semantic reindexing}
\label{sec:residual-restriction-semantic-reindexing}

\subsection{Residual slice comparison}
\label{subsec:residual-slice-comparison}

Let
\[
    F:\mathbb A'\to\mathbb A
\]
be an internal functor.

For a generalized object
\[
    v'\in A'_0,
\]
there is an induced internal functor
\[
    F_{/v'}:
    \mathbb A'/v'
    \to
    \mathbb A/F_0v'
\]
which sends
\[
    r:u'\to v'
\]
to
\[
    F_1r:F_0u'\to F_0v'.
\]

Precomposition gives
\[
    (F_{/v'})^\ast:
    [\mathbb A/F_0v',\mathbb B]
    \to
    [\mathbb A'/v',\mathbb B].
\]

Systems-theoretically, a residual behaviour at \(v'\) describes the future
response of a state to every input continuation entering \(v'\).
Precomposition along \(F_{/v'}\) forgets the global continuations which are
not represented by the local interface.

\subsection{Residual restriction}
\label{subsec:residual-restriction}

\begin{definition}[Residual restriction comparison]
\label{def:residual-restriction-comparison}
The \textbf{residual restriction comparison} is the carrier map
\[
    F^\sharp:
    \overline F^{\,\ast}
    \mathbf{Beh}_{\mathbb A,\mathbb B,0}
    \to
    \mathbf{Beh}_{\mathbb A',\mathbb B,0}
\]
whose component over \(v'\) sends
\[
    H:
    \mathbb A/F_0v'\to\mathbb B
\]
to
\[
    H\circ F_{/v'}:
    \mathbb A'/v'\to\mathbb B.
\]
\end{definition}

The source is the carrier of
\[
    F^\ast
    \mathbf{Beh}^{\mathrm{can}}_{\mathbb A,\mathbb B}.
\]

\begin{theorem}[Strictness of residual restriction]
\label{thm:strictness-residual-restriction}
For every internal functor
\[
    F:\mathbb A'\to\mathbb A,
\]
the map \(F^\sharp\) is a strict semantic homomorphism
\[
    F^\sharp:
    F^\ast
    \mathbf{Beh}^{\mathrm{can}}_{\mathbb A,\mathbb B}
    \to
    \mathbf{Beh}^{\mathrm{can}}_{\mathbb A',\mathbb B}.
\]
Explicitly,
\[
    F^\sharp
    \left(
        \partial_{F_1a'}H
    \right)
    =
    \partial_{a'}
    \left(
        F^\sharp H
    \right),
\]
and
\[
    \omega_{\mathbb A'}
    \left(
        a',
        F^\sharp H
    \right)
    =
    \omega_{\mathbb A}
    \left(
        F_1a',
        H
    \right).
\]
\end{theorem}

\begin{proof}
For
\[
    a':u'\to v',
\]
the residual functors fit into the strictly commuting square
\[
\begin{tikzcd}[column sep=large,row sep=large]
    \mathbb A'/u'
        \arrow[r]
        \arrow[d, "F_{/u'}"']
    &
    \mathbb A'/v'
        \arrow[d, "F_{/v'}"]
    \\
    \mathbb A/F_0u'
        \arrow[r]
    &
    \mathbb A/F_0v'.
\end{tikzcd}
\]
Precomposition around this square gives the derivative equation.

The output comparison is obtained by evaluating the residual functor on the
distinguished arrow represented by the current input. The functor
\(F_{/v'}\) sends the distinguished arrow represented by \(a'\) to that
represented by \(F_1a'\). Hence the output equation holds.
\end{proof}

\subsection{Coherence of residual restriction}
\label{subsec:coherence-residual-restriction}

For composable input functors
\[
\begin{tikzcd}[column sep=large]
    \mathbb A''
        \arrow[r, "G"]
    &
    \mathbb A'
        \arrow[r, "F"]
    &
    \mathbb A,
\end{tikzcd}
\]
residual precomposition satisfies
\[
    (F\circ G)^\sharp
    =
    G^\sharp
    \circ
    \overline G^{\,\ast}F^\sharp
\]
after the canonical pullback comparison.

\begin{proposition}[Pseudonaturality of residual restriction]
\label{prop:pseudonaturality-residual-restriction}
The family \(F^\sharp\) satisfies
\[
    1_{\mathbb A}^\sharp
    =
    1,
\]
and
\[
    (F\circ G)^\sharp
    =
    G^\sharp
    \circ
    \overline G^{\,\ast}F^\sharp
\]
up to canonical pullback coherence.
\end{proposition}

\begin{proof}
Both composites act on a residual functor by precomposition with
\[
    \mathbb A''/v''
    \to
    \mathbb A'/G_0v''
    \to
    \mathbb A/F_0G_0v''.
\]
\end{proof}

\subsection{Semantic direct reindexing}
\label{subsec:semantic-direct-reindexing}

Let
\[
    V
    \hookrightarrow
    \mathbf{Beh}_{\mathbb A,\mathbb B,0}
\]
be a behavioural predicate.

Ordinary base pullback gives
\[
    \overline F^{\,\ast}V
    \hookrightarrow
    \overline F^{\,\ast}
    \mathbf{Beh}_{\mathbb A,\mathbb B,0},
\]
but its codomain is not the canonical behaviour object over
\(\mathbb A'\).

\begin{definition}[Semantic direct reindexing]
\label{def:semantic-direct-reindexing}
Define
\[
    F^\bullet V
    \hookrightarrow
    \mathbf{Beh}_{\mathbb A',\mathbb B,0}
\]
to be the regular image of
\[
\begin{tikzcd}[column sep=large]
    \overline F^{\,\ast}V
        \arrow[r, tail]
    &
    \overline F^{\,\ast}
    \mathbf{Beh}_{\mathbb A,\mathbb B,0}
        \arrow[r, "F^\sharp"]
    &
    \mathbf{Beh}_{\mathbb A',\mathbb B,0}.
\end{tikzcd}
\]
\end{definition}

Thus \(F^\bullet V\) consists of the local residual behaviours which actually
arise by restricting members of \(V\).

\begin{proposition}[Basic properties of semantic direct reindexing]
\label{prop:basic-properties-semantic-direct-reindexing}
Semantic direct reindexing:
\begin{enumerate}
    \item is monotone;
    \item preserves arbitrary joins;
    \item satisfies
    \[
        1_{\mathbb A}^\bullet
        =
        1;
    \]
    \item satisfies a canonical composition isomorphism
    \[
        G^\bullet F^\bullet V
        \cong
        (F\circ G)^\bullet V.
    \]
\end{enumerate}
\end{proposition}

\begin{proof}
Monotonicity and join preservation follow from pullback and regular-image
formation in a Grothendieck topos. The identity comparison follows from
\(1^\sharp=1\).

For composition, pullback preserves regular images, and the image of a
composite is unchanged by replacing the intermediate domain by its regular
image. The result then follows from
Proposition~\ref{prop:pseudonaturality-residual-restriction}.
\end{proof}

\subsection{Naturality of terminal semantics}
\label{subsec:naturality-terminal-semantics-input}

Let
\[
    \beta_\Phi:
    P\to
    \mathbf{Beh}_{\mathbb A,\mathbb B,0}
\]
be the terminal semantic map of
\[
    \Phi:\mathbb A\rightsquigarrow\mathbb B.
\]

\begin{proposition}[Naturality of residual semantics]
\label{prop:naturality-residual-semantics}
The square
\[
\begin{tikzcd}[column sep=large,row sep=large]
    F^\ast P
        \arrow[r, "\overline F^{\,\ast}\beta_\Phi"]
        \arrow[d, "\beta_{F^\ast\Phi}"']
    &
    \overline F^{\,\ast}
    \mathbf{Beh}_{\mathbb A,\mathbb B,0}
        \arrow[d, "F^\sharp"]
    \\
    \mathbf{Beh}_{\mathbb A',\mathbb B,0}
        \arrow[r, equal]
    &
    \mathbf{Beh}_{\mathbb A',\mathbb B,0}
\end{tikzcd}
\]
commutes.
\end{proposition}

\begin{proof}
For a pulled-back state \((v',x)\), the upper-right composite computes the
global residual behaviour of \(x\) and restricts it along
\[
    F_{/v'}.
\]
The lower-left composite computes the residual behaviour of the restricted
system from \((v',x)\). Both assign to each local residual input path its
image under \(F\), followed by the global residual semantics of \(x\).
\end{proof}

\begin{theorem}[Behavioural abstraction under arbitrary restriction]
\label{thm:behavioural-abstraction-arbitrary-restriction}
For every internal functor
\[
    F:\mathbb A'\to\mathbb A
\]
and every system
\[
    \Phi:\mathbb A\rightsquigarrow\mathbb B,
\]
there is a canonical equality of subobjects
\[
    \operatorname{Beh}_{\mathbb A'}(F^\ast\Phi)
    =
    F^\bullet
    \operatorname{Beh}_{\mathbb A}(\Phi).
\]
\end{theorem}

\begin{proof}
The semantic image of \(F^\ast\Phi\) is the regular image of
\[
    \beta_{F^\ast\Phi}
    =
    F^\sharp
    \circ
    \overline F^{\,\ast}\beta_\Phi.
\]
The regular image of
\(\overline F^{\,\ast}\beta_\Phi\) is
\[
    \overline F^{\,\ast}
    \operatorname{Beh}_{\mathbb A}(\Phi)
\]
because regular images are stable under pullback. Taking the image after
\(F^\sharp\) gives
\[
    F^\bullet
    \operatorname{Beh}_{\mathbb A}(\Phi).
\]
\end{proof}

\section{Residual cover classes}
\label{sec:residual-cover-classes}

\subsection{Behaviour-exact input functors}
\label{subsec:behaviour-exact-input-functors}

\begin{definition}[Behaviour-exact input functor]
\label{def:behaviour-exact-input-functor}
An internal functor
\[
    F:\mathbb A'\to\mathbb A
\]
is \textbf{behaviour-exact relative to \(\mathbb B\)} if
\[
    F^\sharp:
    F^\ast
    \mathbf{Beh}^{\mathrm{can}}_{\mathbb A,\mathbb B}
    \xrightarrow{\cong}
    \mathbf{Beh}^{\mathrm{can}}_{\mathbb A',\mathbb B}
\]
is an isomorphism of strict Mealy systems.
\end{definition}

Behaviour exactness says that the local interface has exactly the correct
future semantics:
\begin{enumerate}
    \item two global residual behaviours cannot become indistinguishable after
    restriction;
    \item every local residual behaviour extends uniquely to a global one over
    the represented global context.
\end{enumerate}

Thus behaviour exactness is a lossless interface translation at the level of
complete future response.

When \(F\) is behaviour-exact, define residual pullback by transporting
ordinary slice pullback through \(F^\sharp\):
\[
    F^\ast V
    :=
    F^\sharp
    \left(
        \overline F^{\,\ast}V
    \right).
\]
Then
\[
    F^\bullet V
    =
    F^\ast V.
\]

\begin{corollary}[Behavioural abstraction under exact restriction]
\label{cor:behavioural-abstraction-exact-restriction}
If \(F\) is behaviour-exact, then
\[
    F^\ast
    \operatorname{Beh}_{\mathbb A}(\Phi)
    \cong
    \operatorname{Beh}_{\mathbb A'}(F^\ast\Phi).
\]
\end{corollary}

\subsection{Incoming-arrow comparison}
\label{subsec:incoming-arrow-comparison}

Define the object of global arrows entering a locally represented target by
\[
    \widehat A_{1,F}
    :=
    A'_0
    \times_{F_0,A_0,t_A}
    A_1.
\]

A generalized element is a pair
\[
    (v',a)
\]
with
\[
    t_A(a)=F_0(v').
\]

There is a canonical map
\[
    \lambda_F:
    A'_1\to\widehat A_{1,F}
\]
given by
\[
    \lambda_F(a')
    =
    \left(
        t_{A'}a',
        F_1a'
    \right).
\]

\begin{definition}[Derivative-covering input functor]
\label{def:derivative-covering-input-functor}
An internal functor
\[
    F:\mathbb A'\to\mathbb A
\]
is \textbf{derivative-covering} if
\[
    \lambda_F:
    A'_1\twoheadrightarrow\widehat A_{1,F}
\]
is a regular epimorphism.
\end{definition}

Derivative covering says:

\begin{quote}
Whenever the local interface represents a current input context, it also
represents every global input action which may be applied in that context.
\end{quote}

This is the interface condition needed for local safety checks to be complete.
Without it, a local system may be safe under every locally visible input while
omitting a globally possible input that violates the contract.

\subsection{Universal residual exactness}
\label{subsec:universal-residual-exactness}

Behaviour exactness and derivative covering need not be stable under arbitrary
base change. The residual topology therefore uses their universal forms.

\begin{definition}[Universally behaviour-exact input functor]
\label{def:universally-behaviour-exact}
An internal functor \(F\) is \textbf{universally behaviour-exact} if every
pullback of \(F\) in \(\mathsf{Inp}\) is behaviour-exact relative to
\(\mathbb B\).
\end{definition}

\begin{definition}[Universally derivative-covering input functor]
\label{def:universally-derivative-covering}
An internal functor \(F\) is \textbf{universally derivative-covering} if every
pullback of \(F\) in \(\mathsf{Inp}\) is derivative-covering.
\end{definition}

\begin{definition}[Residually effective cover]
\label{def:residually-effective-cover}
An internal functor
\[
    p:\mathbb U\to\mathbb A
\]
is a \textbf{residually effective cover relative to \(\mathbb B\)} if:
\begin{enumerate}
    \item \(p\) is a \(2\)-nerve cover;
    \item \(p\) is universally behaviour-exact;
    \item \(p\) is universally derivative-covering.
\end{enumerate}
\end{definition}

A \textbf{nervewise residually effective cover} is a residually effective
cover which is also nervewise covering.

\subsection{Stability of residual covers}
\label{subsec:stability-residual-covers}

\begin{lemma}[Composition of derivative-covering functors]
\label{lem:composition-derivative-covering}
If
\[
\begin{tikzcd}[column sep=large]
    \mathbb A''
        \arrow[r, "G"]
    &
    \mathbb A'
        \arrow[r, "F"]
    &
    \mathbb A
\end{tikzcd}
\]
are derivative-covering, then \(F\circ G\) is derivative-covering.
\end{lemma}

\begin{proof}
The incoming-arrow comparison for the composite factors as
\[
\begin{tikzcd}[column sep=large]
    A''_1
        \arrow[r, "\lambda_G", two heads]
    &
    A''_0\times_{G_0,A'_0,t_{A'}}A'_1
        \arrow[r, two heads]
    &
    A''_0\times_{F_0G_0,A_0,t_A}A_1.
\end{tikzcd}
\]
The second map is the pullback of \(\lambda_F\) along \(G_0\). Both maps are
regular epimorphisms.
\end{proof}

\begin{lemma}[Composition of behaviour-exact functors]
\label{lem:composition-behaviour-exact}
If \(F\) and \(G\) are behaviour-exact, then \(F\circ G\) is
behaviour-exact.
\end{lemma}

\begin{proof}
By
Proposition~\ref{prop:pseudonaturality-residual-restriction},
the residual restriction map of the composite is the composite of the
pullback of \(F^\sharp\) with \(G^\sharp\). A composite of isomorphisms is an
isomorphism.
\end{proof}

\begin{proposition}[Stability of residually effective covers]
\label{prop:stability-residually-effective-covers}
Residually effective covers:
\begin{enumerate}
    \item contain identities;
    \item are stable under pullback;
    \item are closed under composition.
\end{enumerate}
The same statements hold for nervewise residually effective covers.
\end{proposition}

\begin{proof}
Pullback stability follows from the universal clauses and
Proposition~\ref{prop:stability-nerve-covers}.

A pullback of a composite can be factored as a composite of successive
pullbacks. Applying
Lemmas~\ref{lem:composition-derivative-covering} and
\ref{lem:composition-behaviour-exact}
after these pullbacks proves composition stability.
\end{proof}

\subsection{The residual interface topology}
\label{subsec:residual-interface-topology}

\begin{definition}[Residual interface pretopology]
\label{def:residual-interface-pretopology}
Let
\[
    \mathcal K_{\mathrm{res}}
\]
be the smallest Grothendieck pretopology on \(\mathsf{Inp}\) containing:
\begin{enumerate}
    \item all extensive interface covers;
    \item all singleton residually effective covers.
\end{enumerate}
Write
\[
    \tau_{\mathrm{res}}
\]
for the generated Grothendieck topology.
\end{definition}

Proposition~\ref{prop:residual-exactness-coproduct-inclusions} ensures that
the members of the extensive generating families have the residual exactness
properties used below.

\section{Descent of residual behavioural contracts}
\label{sec:descent-residual-behavioural-contracts}

Throughout this section, let
\[
    p:\mathbb U\to\mathbb A
\]
be a residually effective cover.

\subsection{Descent of behavioural predicates}
\label{subsec:descent-behavioural-predicates-rigorous}

Behaviour exactness gives an isomorphism
\[
    p^\sharp:
    \overline p^{\,\ast}
    \mathbf{Beh}_{\mathbb A,\mathbb B,0}
    \xrightarrow{\cong}
    \mathbf{Beh}_{\mathbb U,\mathbb B,0}.
\]

\begin{definition}[Predicate descent structure]
\label{def:predicate-descent-datum}
A \textbf{predicate descent structure} consists of a subobject
\[
    C_{\mathbb U}
    \hookrightarrow
    \mathbf{Beh}_{\mathbb U,\mathbb B,0}
\]
whose two pullbacks to
\(\mathbb U^{\langle2\rangle}\) agree under the canonical behaviour-exact
comparison.
\end{definition}

Write
\[
    \mathsf{Desc}_p
    \left(
        \mathsf{Pred}_{\mathbb B}
    \right)
\]
for the resulting poset.

\begin{proposition}[Descent of behavioural predicates]
\label{prop:descent-behavioural-predicates-rigorous}
The comparison
\[
    \mathsf{Pred}_{\mathbb B}(\mathbb A)
    \longrightarrow
    \mathsf{Desc}_p
    \left(
        \mathsf{Pred}_{\mathbb B}
    \right)
\]
is an isomorphism of posets.
\end{proposition}

\begin{proof}
The map
\[
    \overline p:
    U_0\times B_0
    \twoheadrightarrow
    A_0\times B_0
\]
is an effective descent morphism for objects and subobjects in the
corresponding slices. Behaviour exactness identifies the local canonical
behaviour object with the pullback of the global one. Compatible local
subobjects therefore descend uniquely.
\end{proof}

\subsection{The pulled-back global derivative domain}
\label{subsec:pulled-back-global-derivative-domain}

Write
\[
    \mathbf{Beh}_{\mathbb A,0}
    :=
    \mathbf{Beh}_{\mathbb A,\mathbb B,0}.
\]

The global derivative domain is
\[
    D_{\mathbb A}
    :=
    A_1
    \times_{t_A,A_0,p_{\mathbf{Beh}}}
    \mathbf{Beh}_{\mathbb A,0}.
\]
Write
\[
    d_0:
    D_{\mathbb A}\to
    \mathbf{Beh}_{\mathbb A,0}
\]
for the current-state projection, and
\[
    d_1:
    D_{\mathbb A}\to
    \mathbf{Beh}_{\mathbb A,0}
\]
for the derivative action.

Define
\[
    \widehat D_p
    :=
    \widehat A_{1,p}
    \times_{U_0}
    \overline p^{\,\ast}
    \mathbf{Beh}_{\mathbb A,0}.
\]

Using behaviour exactness, the local derivative domain is
\[
    D_{\mathbb U}
    :=
    U_1
    \times_{t_U,U_0,p_{\mathbf{Beh}}}
    \mathbf{Beh}_{\mathbb U,0}.
\]

There is a canonical comparison
\[
    \Lambda_p:
    D_{\mathbb U}\to\widehat D_p
\]
induced by
\[
    \lambda_p:
    U_1\to\widehat A_{1,p}
\]
and the inverse of \(p^\sharp\).

\begin{lemma}[Derivative-domain covering]
\label{lem:derivative-domain-covering}
The map
\[
    \Lambda_p:
    D_{\mathbb U}\twoheadrightarrow\widehat D_p
\]
is a regular epimorphism.
\end{lemma}

\begin{proof}
The map \(\Lambda_p\) is a pullback of
\[
    \lambda_p:
    U_1\twoheadrightarrow\widehat A_{1,p}
\]
along the pulled-back canonical behaviour carrier. Regular epimorphisms are
stable under pullback.
\end{proof}

\subsection{Derivative invariance as a subobject inclusion}
\label{subsec:derivative-invariance-subobject-inclusion}

Let
\[
    C\hookrightarrow\mathbf{Beh}_{\mathbb A,0}
\]
be a behavioural predicate.

\begin{proposition}[Subobject form of derivative invariance]
\label{prop:subobject-form-derivative-invariance}
The predicate \(C\) is derivative-invariant if and only if
\[
    d_0^\ast C
    \leq
    d_1^\ast C
\]
in
\[
    \operatorname{Sub}(D_{\mathbb A}).
\]
\end{proposition}

\begin{proof}
The subobject \(d_0^\ast C\) consists of current residual states in \(C\)
together with applicable input arrows. The subobject \(d_1^\ast C\) consists
of those pairs whose derivative state lies in \(C\). The inclusion says
exactly that every derivative of a state in \(C\) remains in \(C\).
\end{proof}

\subsection{Locality of derivative invariance}
\label{subsec:locality-derivative-invariance-rigorous}

\begin{theorem}[Derivative invariance is interface-local]
\label{thm:derivative-invariance-interface-local}
A behavioural predicate
\[
    C\hookrightarrow\mathbf{Beh}_{\mathbb A,0}
\]
is derivative-invariant if and only if
\[
    p^\ast C
    \hookrightarrow
    \mathbf{Beh}_{\mathbb U,0}
\]
is derivative-invariant.
\end{theorem}

\begin{proof}
If \(C\) is derivative-invariant, strictness of \(p^\sharp\) implies that
\(p^\ast C\) is derivative-invariant.

Conversely, suppose \(p^\ast C\) is derivative-invariant. Its invariance is
the inclusion
\[
    d_{0,\mathbb U}^\ast(p^\ast C)
    \leq
    d_{1,\mathbb U}^\ast(p^\ast C)
\]
on \(D_{\mathbb U}\).

Under behaviour exactness, this is the pullback along
\[
    \Lambda_p:
    D_{\mathbb U}\twoheadrightarrow\widehat D_p
\]
of the pullback to \(\widehat D_p\) of
\[
    d_0^\ast C
    \leq
    d_1^\ast C.
\]
Pullback along a regular epimorphism reflects subobject inclusion, so the
candidate inclusion holds on \(\widehat D_p\).

The object \(\widehat D_p\) is the pullback of \(D_{\mathbb A}\) along the
regular epimorphism covering current-state anchors. Pullback along this map
again reflects subobject inclusion. Hence
\[
    d_0^\ast C
    \leq
    d_1^\ast C
\]
holds globally.
\end{proof}

\begin{remark}[Safety interpretation]
\label{rem:interface-local-safety-interpretation}
The forward implication says that a globally safe residual specification
remains safe after restricting the input vocabulary.

The converse is more substantial. It says that local safety is complete
because:
\begin{enumerate}
    \item behaviour exactness ensures that the local residual states are
    exactly the restrictions of global residual states;
    \item derivative covering ensures that no globally possible input at a
    represented context has been omitted from the local safety test;
    \item object coverage ensures that every global current context is
    represented somewhere locally.
\end{enumerate}
\end{remark}

\subsection{Contract descent}
\label{subsec:contract-descent-rigorous}

\begin{definition}[Contract descent structure]
\label{def:contract-descent-datum-rigorous}
A \textbf{contract descent structure} is a predicate descent structure
\[
    V_{\mathbb U}
    \hookrightarrow
    \mathbf{Beh}_{\mathbb U,0}
\]
such that \(V_{\mathbb U}\) is derivative-invariant.
\end{definition}

Write
\[
    \mathsf{Desc}_p
    \left(
        \mathsf{Ctr}_{\mathbb B}
    \right)
\]
for the resulting poset.

\begin{theorem}[Effective descent of residual contracts]
\label{thm:effective-descent-residual-contracts}
For every residually effective cover
\[
    p:\mathbb U\to\mathbb A,
\]
the comparison
\[
    \mathsf{Ctr}_{\mathbb B}(\mathbb A)
    \longrightarrow
    \mathsf{Desc}_p
    \left(
        \mathsf{Ctr}_{\mathbb B}
    \right)
\]
is an isomorphism of posets.
\end{theorem}

\begin{proof}
By
Proposition~\ref{prop:descent-behavioural-predicates-rigorous},
every compatible local predicate descends uniquely to a global predicate.
By
Theorem~\ref{thm:derivative-invariance-interface-local},
the descended predicate is derivative-invariant exactly when its local
pullback is derivative-invariant.
\end{proof}

\begin{remark}[Specifications versus implementations]
\label{rem:specifications-versus-implementations}
Contract descent reconstructs a global set of permitted residual behaviours.
It does not reconstruct a chosen implementation unless the local systems also
carry compatible state and transition identifications.

Thus:
\[
    \text{compatible local specifications}
    \centernot\Longrightarrow
    \text{compatible chosen implementations}.
\]
Several locally chosen controllers may satisfy contracts that glue perfectly
while using incompatible hidden states or incompatible overlap dynamics.
\end{remark}

\begin{corollary}[Stack of residual contracts]
\label{cor:stack-residual-contracts}
The indexed poset
\[
    \mathsf{Ctr}_{\mathbb B}
\]
is a stack for
\[
    \tau_{\mathrm{res}}.
\]
Equivalently, it is a sheaf of posets on the residual interface site.
\end{corollary}

\begin{proof}
Extensive descent is
Theorem~\ref{thm:extensive-descent-residual-contracts}.
Singleton residual descent is
Theorem~\ref{thm:effective-descent-residual-contracts}.
The generating cover classes are pullback-stable, and descent composes by
Lemma~\ref{lem:composition-effective-descent}.
\end{proof}

\section{Derivative closure and canonical realization under restriction}
\label{sec:derivative-closure-canonical-realization-restriction}

Throughout this section, let
\[
    F:\mathbb A'\to\mathbb A
\]
be behaviour-exact and derivative-covering.

\subsection{Generated derivative closure}
\label{subsec:generated-derivative-closure-restriction}

Let
\[
    C\hookrightarrow\mathbf{Beh}_{\mathbb A,0}.
\]

Define
\[
    D_{\mathbb A}(C)
    :=
    D_{\mathbb A}
    \times_{d_0,\mathbf{Beh}_{\mathbb A,0}}
    C.
\]
Generated derivative closure is the regular image
\[
    \operatorname{Gen}_{\partial,\mathbb A}(C)
    =
    \operatorname{Im}
    \left(
        D_{\mathbb A}(C)
        \xrightarrow{d_1}
        \mathbf{Beh}_{\mathbb A,0}
    \right).
\]

Because identity arrows are included, this image contains \(C\). Because
arbitrary input arrows include composites, it is derivative-invariant.

\begin{lemma}[Images under regular-epimorphic refinement]
\label{lem:images-regular-epimorphic-refinement}
Let
\[
    E\twoheadrightarrow D\xrightarrow{f}X
\]
be a regular epimorphism followed by an arbitrary map. Then
\[
    \operatorname{Im}(f)
    =
    \operatorname{Im}(f\circ e).
\]
\end{lemma}

\begin{proof}
Factor \(f\) as a regular epimorphism followed by a monomorphism. Since the
composite of regular epimorphisms is regular epic, the same monomorphism is an
image factor for \(f\circ e\).
\end{proof}

\begin{theorem}[Generated closure commutes with residual restriction]
\label{thm:generated-closure-commutes-restriction}
There is a canonical isomorphism
\[
    F^\ast
    \operatorname{Gen}_{\partial,\mathbb A}(C)
    \cong
    \operatorname{Gen}_{\partial,\mathbb A'}
    (F^\ast C).
\]
\end{theorem}

\begin{proof}
Pullback stability of regular images identifies
\[
    F^\ast
    \operatorname{Gen}_{\partial,\mathbb A}(C)
\]
with the image of the pulled-back global derivative action on the pulled-back
predicate.

The pulled-back global incoming-arrow domain is
\(\widehat D_F\). The local derivative domain maps to it by the regular
epimorphism
\[
    \Lambda_F:
    D_{\mathbb A'}\twoheadrightarrow\widehat D_F.
\]
Strictness of \(F^\sharp\) identifies the local derivative map with the
pulled-back global derivative map after precomposition by \(\Lambda_F\).

By
Lemma~\ref{lem:images-regular-epimorphic-refinement},
this precomposition does not change the image. The resulting image is
\[
    \operatorname{Gen}_{\partial,\mathbb A'}
    (F^\ast C).
\]
\end{proof}

\subsection{Safety cores}
\label{subsec:safety-cores-restriction}

The safety core is
\[
    \operatorname{Safe}_{\partial,\mathbb A}(C)
    =
    \forall_{d_0}
    \left(
        d_1^\ast C
    \right).
\]

\begin{lemma}[Universal quantification under an epimorphic refinement]
\label{lem:universal-quantification-epimorphic-refinement}
Let
\[
    E\xrightarrow{e}D\xrightarrow{f}X
\]
be maps in a Grothendieck topos, with \(e\) a regular epimorphism. For every
subobject
\[
    S\hookrightarrow D,
\]
there is a canonical equality
\[
    \forall_f(S)
    =
    \forall_{fe}(e^\ast S).
\]
\end{lemma}

\begin{proof}
For every subobject \(T\hookrightarrow X\),
\[
\begin{aligned}
    T\leq\forall_{fe}(e^\ast S)
    &\iff
    e^\ast f^\ast T\leq e^\ast S
    \\
    &\iff
    f^\ast T\leq S
    \\
    &\iff
    T\leq\forall_f(S).
\end{aligned}
\]
The middle equivalence uses reflection of subobject inclusion by pullback
along the regular epimorphism \(e\). Hence the two right adjoints define the
same subobject.
\end{proof}

\begin{theorem}[Safety core commutes with residual restriction]
\label{thm:safety-core-commutes-restriction}
There is a canonical isomorphism
\[
    F^\ast
    \operatorname{Safe}_{\partial,\mathbb A}(C)
    \cong
    \operatorname{Safe}_{\partial,\mathbb A'}
    (F^\ast C).
\]
\end{theorem}

\begin{proof}
The square comparing the pulled-back global derivative domain
\(\widehat D_F\) with the global derivative domain is a pullback.
Beck--Chevalley for dependent products gives
\[
    F^\ast
    \forall_{d_0}
    \left(
        d_1^\ast C
    \right)
    \cong
    \forall_{\widehat d_0}
    \left(
        \widehat d_1^\ast F^\ast C
    \right).
\]

The local derivative domain maps by the regular epimorphism
\[
    \Lambda_F:
    D_{\mathbb A'}\twoheadrightarrow\widehat D_F.
\]
Strictness identifies the pullback of
\(\widehat d_1^\ast F^\ast C\) with
\[
    d_{1,\mathbb A'}^\ast(F^\ast C).
\]
Applying
Lemma~\ref{lem:universal-quantification-epimorphic-refinement}
gives
\[
    \forall_{\widehat d_0}
    \left(
        \widehat d_1^\ast F^\ast C
    \right)
    =
    \forall_{d_{0,\mathbb A'}}
    \left(
        d_{1,\mathbb A'}^\ast F^\ast C
    \right).
\]
\end{proof}

\subsection{The indexed derivative adjoint triple}
\label{subsec:indexed-derivative-adjoint-triple}

For each input category \(\mathbb A\), recall
\[
    \operatorname{Gen}_{\partial,\mathbb A}
    \dashv
    I_{\partial,\mathbb A}
    \dashv
    \operatorname{Safe}_{\partial,\mathbb A}.
\]

\begin{corollary}[Indexed derivative adjoint triple]
\label{cor:indexed-derivative-adjoint-triple}
Over the category whose morphisms are behaviour-exact,
derivative-covering input functors, the fibrewise adjunctions
\[
    \operatorname{Gen}_{\partial}
    \dashv
    I_{\partial}
    \dashv
    \operatorname{Safe}_{\partial}
\]
assemble into indexed adjunctions. Their Beck--Chevalley comparisons are
invertible.
\end{corollary}

\begin{proof}
The invertible comparisons for generated closure and safety are
Theorems~\ref{thm:generated-closure-commutes-restriction} and
\ref{thm:safety-core-commutes-restriction}. The inclusion functors commute
strictly with pullback. Naturality of units and counits follows from
uniqueness of adjoint mates.
\end{proof}

\subsection{Restriction of canonical realizations}
\label{subsec:restriction-canonical-realizations}

Let
\[
    V\hookrightarrow\mathbf{Beh}_{\mathbb A,0}
\]
be derivative-invariant. Its canonical realization
\[
    \iota_{\mathbb A}(V)
\]
is the restriction of the canonical behaviour Mealy morphism to \(V\).

\begin{proposition}[Restriction of canonical realizations]
\label{prop:restriction-canonical-realizations}
There is a canonical strict isomorphism
\[
    F^\ast
    \iota_{\mathbb A}(V)
    \cong
    \iota_{\mathbb A'}
    (F^\ast V).
\]
\end{proposition}

\begin{proof}
Behaviour exactness identifies the pulled-back canonical behaviour machine
with the local canonical behaviour machine. Under this identification,
\(F^\ast V\) is derivative-invariant by strictness of \(F^\sharp\), and the
restricted transition and output maps coincide.
\end{proof}

\subsection{Indexed abstraction and realization}
\label{subsec:indexed-abstraction-realization}

For every \(\mathbb A\), recall the fibrewise adjunction
\[
    \operatorname{Beh}_{\mathbb A}
    \dashv
    \iota_{\mathbb A}.
\]

\begin{theorem}[Indexed abstraction--realization adjunction]
\label{thm:indexed-abstraction-realization-adjunction}
Over the category of residually effective input functors, the fibrewise
adjunctions
\[
    \operatorname{Beh}_{\mathbb A}
    \dashv
    \iota_{\mathbb A}
\]
assemble into an indexed adjunction between:
\begin{enumerate}
    \item the indexed category of strict Mealy systems;
    \item the indexed poset of residual behavioural contracts.
\end{enumerate}
\end{theorem}

\begin{proof}
The left-adjoint comparison is
\[
    F^\ast
    \operatorname{Beh}_{\mathbb A}(\Phi)
    \cong
    \operatorname{Beh}_{\mathbb A'}(F^\ast\Phi)
\]
from
Corollary~\ref{cor:behavioural-abstraction-exact-restriction}.
The right-adjoint comparison is
Proposition~\ref{prop:restriction-canonical-realizations}.
They are mates under the fibrewise adjunctions and satisfy the coherence
inherited from input restriction and residual restriction.
\end{proof}

\section{Locality of satisfaction and minimization}
\label{sec:locality-satisfaction-minimization}

Throughout this section, let
\[
    \{F_i:\mathbb A_i\to\mathbb A\}_{i\in I}
\]
be a covering family in
\(\mathcal K_{\mathrm{res}}\).

Every member of such a family is obtained from extensive inclusions and
residually effective singleton covers by pullback and composition.
Consequently, each \(F_i\) is behaviour-exact and derivative-covering.

\subsection{Joint reflection of subobject inclusion}
\label{subsec:joint-reflection-subobject-inclusion}

\begin{lemma}[Local reflection of subobject inclusion]
\label{lem:local-reflection-subobject-inclusion}
Let
\[
    S,T
    \hookrightarrow
    \mathbf{Beh}_{\mathbb A,0}
\]
be subobjects. If
\[
    F_i^\ast S
    \leq
    F_i^\ast T
\]
for every \(i\), then
\[
    S\leq T.
\]
\end{lemma}

\begin{proof}
For a singleton residually effective cover, pullback along the regular
epimorphism on the behaviour base reflects subobject inclusion.

For an extensive cover, subobjects decompose componentwise.

The reflection property is preserved under pullback and composition of
covering families, so it holds for every family in the generated
pretopology.
\end{proof}

\subsection{Locality of contract satisfaction}
\label{subsec:locality-contract-satisfaction-rigorous}

Let
\[
    \Phi:\mathbb A\rightsquigarrow\mathbb B
\]
and let
\[
    V\hookrightarrow\mathbf{Beh}_{\mathbb A,0}
\]
be a residual contract.

Recall that
\[
    \Phi\models V
\]
means
\[
    \operatorname{Beh}_{\mathbb A}(\Phi)\leq V.
\]

\begin{theorem}[Locality of contract satisfaction]
\label{thm:locality-contract-satisfaction-rigorous}
The following are equivalent:
\begin{enumerate}
    \item
    \[
        \Phi\models V;
    \]
    \item for every \(i\in I\),
    \[
        F_i^\ast\Phi
        \models
        F_i^\ast V.
    \]
\end{enumerate}
\end{theorem}

\begin{proof}
If
\[
    \operatorname{Beh}_{\mathbb A}(\Phi)\leq V,
\]
then pullback gives
\[
    F_i^\ast
    \operatorname{Beh}_{\mathbb A}(\Phi)
    \leq
    F_i^\ast V.
\]
Behaviour exactness identifies the left side with
\[
    \operatorname{Beh}_{\mathbb A_i}(F_i^\ast\Phi).
\]

Conversely, local satisfaction gives
\[
    F_i^\ast
    \operatorname{Beh}_{\mathbb A}(\Phi)
    \leq
    F_i^\ast V
\]
for every \(i\). Apply
Lemma~\ref{lem:local-reflection-subobject-inclusion}.
\end{proof}

\begin{remark}[Interpretation as complete local verification]
\label{rem:complete-local-verification}
The theorem says that a global behavioural contract may be checked patch by
patch without losing any global case.

This completeness depends on the cover conditions. Behaviour exactness
prevents distinct global futures from becoming indistinguishable locally, and
derivative covering prevents a globally possible input from being omitted
from every local check.
\end{remark}

\subsection{Locality of behavioural separatedness}
\label{subsec:locality-behavioural-separatedness}

Recall that \(\Phi\) is behaviourally separated if
\[
    \beta_\Phi:
    P\to
    \mathbf{Beh}_{\mathbb A,0}
\]
is monic.

\begin{theorem}[Separatedness is interface-local]
\label{thm:separatedness-interface-local}
The system \(\Phi\) is behaviourally separated if and only if every local
restriction
\[
    F_i^\ast\Phi
\]
is behaviourally separated.
\end{theorem}

\begin{proof}
Pullback preserves monomorphisms, so global separatedness implies local
separatedness.

Conversely,
Proposition~\ref{prop:naturality-residual-semantics}
and behaviour exactness identify the local semantic map with the pullback of
\(\beta_\Phi\), followed by an isomorphism. Hence every local pullback of
\(\beta_\Phi\) is monic.

Monomorphisms are local for regular-epimorphic and extensive covers. Therefore
\(\beta_\Phi\) is monic.
\end{proof}

\subsection{Locality of semantic minimization}
\label{subsec:locality-semantic-minimization}

Let
\[
    \eta_\Phi:
    \Phi\twoheadrightarrow
    \operatorname{Beh}_{\mathbb A}(\Phi)
\]
be the semantic quotient.

\begin{theorem}[Semantic minimization commutes with local restriction]
\label{thm:semantic-minimization-local-restriction}
For every residually effective input functor
\[
    F:\mathbb A'\to\mathbb A,
\]
there is a canonical isomorphism of semantic quotient factorizations
\[
\begin{tikzcd}[column sep=large,row sep=large]
    F^\ast\Phi
        \arrow[r, two heads, "F^\ast\eta_\Phi"]
        \arrow[dr, two heads, "\eta_{F^\ast\Phi}"']
    &
    F^\ast
    \operatorname{Beh}_{\mathbb A}(\Phi)
        \arrow[d, "\cong"]
    \\
    {}
    &
    \operatorname{Beh}_{\mathbb A'}
    (F^\ast\Phi).
\end{tikzcd}
\]
\end{theorem}

\begin{proof}
The semantic quotient is the regular-image factorization of
\(\beta_\Phi\). Regular images are stable under pullback. Naturality of
residual semantics and behaviour exactness identify the pulled-back
factorization with the regular-image factorization of
\(\beta_{F^\ast\Phi}\).
\end{proof}

\begin{remark}[Automata-theoretic interpretation]
\label{rem:local-minimization-interpretation}
Under behaviour exactness, minimizing globally and then restricting the input
interface gives the same local automaton as restricting first and minimizing
locally.

Without behaviour exactness, local restriction can hide distinguishing input
experiments. Two globally distinct states may then become locally equivalent,
and local minimization may identify more states than restricted global
minimization.
\end{remark}

\begin{corollary}[Local reconstruction of minimal separated quotients]
\label{cor:local-reconstruction-minimal-separated-quotients}
The minimal separated quotient of a global system is reconstructed by
effective descent from the minimal separated quotients of its compatible local
restrictions.
\end{corollary}

\subsection{Locality of generated residual theories}
\label{subsec:locality-generated-residual-theories}

Let
\[
    K
    \hookrightarrow
    \mathbf{Beh}_{\mathbb A,0}
\]
be a behavioural specification. Its derivative-generated theory is
\[
    \operatorname{Der}_0(K)
    :=
    \operatorname{Gen}_{\partial,\mathbb A}(K).
\]

\begin{proposition}[Generated theories commute with local restriction]
\label{prop:generated-theories-local-restriction}
For every residually effective
\[
    F:\mathbb A'\to\mathbb A,
\]
there is a canonical isomorphism
\[
    F^\ast
    \operatorname{Der}_0(K)
    \cong
    \operatorname{Der}_0
    (F^\ast K).
\]
\end{proposition}

\begin{proof}
This is
Theorem~\ref{thm:generated-closure-commutes-restriction}
applied to \(K\).
\end{proof}

\subsection{Locality of residual reachability}
\label{subsec:locality-reachability}

Suppose that \(\Phi\) realizes \(K\). Recall that \(\Phi\) is residually
generated by \(K\) when
\[
    \operatorname{Beh}_{\mathbb A}(\Phi)
    =
    \operatorname{Der}_0(K).
\]

\begin{theorem}[Residual reachability is interface-local]
\label{thm:reachability-interface-local}
The realization \(\Phi\) is residually generated by \(K\) if and only if
every local realization
\[
    F_i^\ast\Phi
\]
is residually generated by
\[
    F_i^\ast K.
\]
\end{theorem}

\begin{proof}
Global equality pulls back to local equality by
Corollary~\ref{cor:behavioural-abstraction-exact-restriction}
and
Proposition~\ref{prop:generated-theories-local-restriction}.

Conversely, local equality gives
\[
    F_i^\ast
    \operatorname{Beh}_{\mathbb A}(\Phi)
    =
    F_i^\ast
    \operatorname{Der}_0(K)
\]
for every \(i\). The cover jointly reflects equality of subobjects.
\end{proof}

\begin{corollary}[Local Myhill--Nerode principle]
\label{cor:local-myhill-nerode-principle}
A realization is residually generated and behaviourally separated if and only
if every restriction to a residual interface cover is residually generated
and behaviourally separated. Consequently, minimal residual realizations may
be verified and reconstructed interface-locally.
\end{corollary}

\section{Trajectory base change}
\label{sec:trajectory-base-change}

This section imports the temporal execution semantics of
Chapter~\ref{ch:temporal-trajectories-segal-execution}.

\subsection{Restriction of execution categories}
\label{subsec:restriction-execution-categories}

Let
\[
    F:\mathbb A'\to\mathbb A
\]
and
\[
    \Phi=(P,\delta_P,\omega_P):
    \mathbb A\rightsquigarrow\mathbb B.
\]

The restricted carrier is
\[
    P_F=A'_0\times_{A_0}P.
\]

Define
\[
    E_{F,\Phi}:
    \operatorname{Exec}(F^\ast\Phi)
    \to
    \operatorname{Exec}(\Phi)
\]
by
\[
    E_{F,\Phi}(v',x)=x
\]
on objects and
\[
    E_{F,\Phi}(a',(v',x))
    =
    (F_1a',x)
\]
on arrows.

\begin{proposition}[Execution restriction functor]
\label{prop:execution-restriction-functor-rigorous}
The displayed maps define an internal functor. Moreover,
\[
    I_\Phi E_{F,\Phi}
    =
    F^{\mathrm{op}}I_{F^\ast\Phi},
\]
and
\[
    O_\Phi E_{F,\Phi}
    =
    O_{F^\ast\Phi}.
\]
\end{proposition}

\begin{proof}
Source preservation is immediate. Target preservation follows from
\[
    \delta_{F^\ast\Phi}(a',(v',x))
    =
    \left(
        s_{A'}a',
        \delta_\Phi(F_1a',x)
    \right).
\]
Identity and composition preservation follow from functoriality of \(F\).
The input-labelling equation is immediate, and the output-labelling equation
follows from
\[
    \omega_{F^\ast\Phi}(a',(v',x))
    =
    \omega_\Phi(F_1a',x).
\]
\end{proof}

\subsection{Cartesian description of execution restriction}
\label{subsec:cartesian-description-execution-restriction}

\begin{theorem}[Execution-category base change]
\label{thm:execution-category-base-change}
There is a canonical pullback isomorphism of internal categories
\[
\begin{tikzcd}[column sep=large,row sep=large]
    \operatorname{Exec}(F^\ast\Phi)
        \arrow[r, "\cong"]
        \arrow[d]
    &
    \operatorname{Exec}(\Phi)
        \arrow[d, "I_\Phi"]
    \\
    \mathbb A'^{\mathrm{op}}
        \arrow[r, "F^{\mathrm{op}}"']
    &
    \mathbb A^{\mathrm{op}}.
\end{tikzcd}
\]
Equivalently,
\[
    \operatorname{Exec}(F^\ast\Phi)
    \cong
    \mathbb A'^{\mathrm{op}}
    \times_{\mathbb A^{\mathrm{op}}}
    \operatorname{Exec}(\Phi).
\]
\end{theorem}

\begin{proof}
On objects, the pullback is
\[
    A'_0\times_{A_0}P,
\]
which is the carrier of \(F^\ast\Phi\).

On arrows, the pullback is
\[
    A'_1
    \times_{A_1}
    \left(
        A_1\times_{A_0}P
    \right),
\]
canonically isomorphic to
\[
    A'_1\times_{A'_0}
    \left(
        A'_0\times_{A_0}P
    \right).
\]
This is the one-step execution object of \(F^\ast\Phi\).

The source, target, identity, and composition maps agree by the defining
formulas for input restriction.
\end{proof}

\subsection{Trajectory base change}
\label{subsec:trajectory-base-change-theorem}

Define
\[
    F^\sharp_{\mathrm{traj}}\mathscr T_\Phi
\]
to be the pullback
\[
\begin{tikzcd}[column sep=large,row sep=large]
    F^\sharp_{\mathrm{traj}}\mathscr T_\Phi
        \arrow[r]
        \arrow[d]
    &
    \mathscr T_\Phi
        \arrow[d, "\operatorname{Path}(I_\Phi)"]
    \\
    \operatorname{Path}(\mathbb A'^{\mathrm{op}})
        \arrow[r, "\operatorname{Path}(F^{\mathrm{op}})"']
    &
    \operatorname{Path}(\mathbb A^{\mathrm{op}}).
\end{tikzcd}
\]

\begin{theorem}[Trajectory base-change theorem]
\label{thm:trajectory-base-change}
There is a canonical isomorphism of labelled trajectory sheaves
\[
    \mathscr T_{F^\ast\Phi}
    \cong
    F^\sharp_{\mathrm{traj}}\mathscr T_\Phi.
\]
At temporal extent \(n\),
\[
    \mathscr T_{F^\ast\Phi}(n)
    \cong
    N(F^{\mathrm{op}})_n^\ast
    \mathscr T_\Phi(n).
\]
\end{theorem}

\begin{proof}
Apply the finite-limit-preserving trajectory functor
\[
    \operatorname{Path}
\]
to the pullback isomorphism of
Theorem~\ref{thm:execution-category-base-change}.
Evaluation at \([n]\) gives the degreewise formula.
\end{proof}

\begin{remark}[Operational meaning of trajectory base change]
\label{rem:operational-trajectory-base-change}
A local finite run is exactly a global finite run whose input path has been
presented through the local interface.

The theorem does not merely compare endpoints. It identifies the complete
state path and output-comparison path after translating each local input step
through \(F\).
\end{remark}

\subsection{Pseudonaturality of trajectory semantics}
\label{subsec:pseudonaturality-trajectory-semantics}

\begin{proposition}[Pseudonaturality of trajectory base change]
\label{prop:pseudonaturality-trajectory-base-change}
The isomorphisms of
Theorem~\ref{thm:trajectory-base-change}
are compatible with:
\begin{enumerate}
    \item identity input functors;
    \item composition of input functors;
    \item strict semantic homomorphisms.
\end{enumerate}
\end{proposition}

\begin{proof}
Each comparison is induced by the canonical pullback comparison for execution
categories. Identity and composition coherence follow from pullback coherence.
Naturality in strict homomorphisms follows from naturality of
\(\operatorname{Exec}(f)\).
\end{proof}

\begin{remark}[Indexed temporal execution]
\label{rem:indexed-temporal-semantics}
The execution-trajectory construction is compatible with changing the input
interface. It therefore forms a pseudonatural family of labelled trajectory
semantics over \(\mathsf{Inp}\).
\end{remark}

\section{Descent of finite trajectories}
\label{sec:descent-finite-trajectories}

Throughout this section, let
\[
    p:\mathbb U\to\mathbb A
\]
be a nervewise cover.

For trajectory labels it is convenient to write
\[
    A_n^\partial
    :=
    N(\mathbb A^{\mathrm{op}})_n,
    \qquad
    U_n^\partial
    :=
    N(\mathbb U^{\mathrm{op}})_n,
\]
and
\[
    p_n^\partial:
    U_n^\partial\to A_n^\partial.
\]
Reversal of a composable string gives canonical isomorphisms
\[
    A_n^\partial\cong A_n,
    \qquad
    U_n^\partial\cong U_n.
\]
Hence every \(p_n^\partial\) is a regular epimorphism.

\subsection{Degreewise trajectory descent}
\label{subsec:degreewise-trajectory-descent}

Let
\[
    \Phi:\mathbb A\rightsquigarrow\mathbb B.
\]
By
Theorem~\ref{thm:trajectory-base-change},
\[
    \mathscr T_{p^\ast\Phi}(n)
    \cong
    (p_n^\partial)^\ast
    \mathscr T_\Phi(n),
\]
where the pullback is taken over the input-path label map.

\begin{proposition}[Effective descent of \(n\)-step executions]
\label{prop:effective-descent-n-step-executions}
For every \(n\geq0\), the object
\[
    \mathscr T_\Phi(n)
    \to
    A_n^\partial
\]
is reconstructed effectively from
\[
    \mathscr T_{p^\ast\Phi}(n)
    \to
    U_n^\partial
\]
and its overlap identification along
\[
    p_n^\partial:
    U_n^\partial\twoheadrightarrow A_n^\partial.
\]
\end{proposition}

\begin{proof}
The local \(n\)-step execution object is the pullback of the global one along
\(p_n^\partial\). Since \(p_n^\partial\) is a regular epimorphism, it is an
effective descent morphism.
\end{proof}

\subsection{Descent of interval restrictions}
\label{subsec:descent-interval-restrictions}

Let
\[
    \iota:[k]\to[n]
\]
be a contiguous interval inclusion. It induces input-path restriction maps
\[
    a_\iota^{\mathbb A}:
    A_n^\partial\to A_k^\partial,
\]
and
\[
    a_\iota^{\mathbb U}:
    U_n^\partial\to U_k^\partial.
\]
They fit into the commutative square
\[
\begin{tikzcd}[column sep=large,row sep=large]
    U_n^\partial
        \arrow[r, "p_n^\partial"]
        \arrow[d, "a_\iota^{\mathbb U}"']
    &
    A_n^\partial
        \arrow[d, "a_\iota^{\mathbb A}"]
    \\
    U_k^\partial
        \arrow[r, "p_k^\partial"']
    &
    A_k^\partial.
\end{tikzcd}
\]

The global trajectory restriction
\[
    r_\iota:
    \mathscr T_\Phi(n)
    \to
    \mathscr T_\Phi(k)
\]
is equivalently a morphism in the slice
\(\mathcal E/A_n^\partial\):
\[
    \widetilde r_\iota:
    \mathscr T_\Phi(n)
    \to
    (a_\iota^{\mathbb A})^\ast
    \mathscr T_\Phi(k).
\]

Likewise, the local restriction is a morphism in
\(\mathcal E/U_n^\partial\):
\[
    \widetilde r_{\iota,\mathbb U}:
    \mathscr T_{p^\ast\Phi}(n)
    \to
    (a_\iota^{\mathbb U})^\ast
    \mathscr T_{p^\ast\Phi}(k).
\]

\begin{proposition}[Descent of temporal restrictions]
\label{prop:descent-temporal-restrictions}
The local temporal restriction maps descend uniquely to the global temporal
restriction maps. The descended maps satisfy the identity and composition
laws of
\[
    \mathbf{Int}_{\mathrm{disc}}^{\mathrm{op}}.
\]
\end{proposition}

\begin{proof}
Trajectory base change and the preceding commutative square give a canonical
identification
\[
    (p_n^\partial)^\ast
    (a_\iota^{\mathbb A})^\ast\mathscr T_\Phi(k)
    \cong
    (a_\iota^{\mathbb U})^\ast
    (p_k^\partial)^\ast\mathscr T_\Phi(k).
\]
Thus the local restriction map is compatible descent structure for a morphism
between two objects of
\(\mathcal E/A_n^\partial\).

Since \(p_n^\partial\) is an effective descent morphism, the local map
descends uniquely to
\[
    \widetilde r_\iota.
\]

The identity and composition equations hold after pullback along
\(p_n^\partial\). Since this pullback functor is faithful, the equations hold
globally.
\end{proof}

\subsection{Descent of the trajectory sheaf}
\label{subsec:descent-trajectory-sheaf}

\begin{theorem}[Effective interface descent of temporal trajectories]
\label{thm:effective-interface-descent-temporal-trajectories}
Let
\[
    p:\mathbb U\to\mathbb A
\]
be a nervewise cover. Compatible local temporal execution sheaves arising
from a system descent structure reconstruct uniquely to the temporal
execution sheaf of the descended global system.
\end{theorem}

\begin{proof}
Since a nervewise cover is a \(2\)-nerve cover, the local system first
reconstructs to a global system
\[
    \Phi:\mathbb A\rightsquigarrow\mathbb B
\]
by
Theorem~\ref{thm:effective-descent-strict-mealy-systems}.

At every temporal degree, trajectory base change identifies the local
execution object with the pullback of the corresponding global object.
Degreewise effective descent is
Proposition~\ref{prop:effective-descent-n-step-executions}.
The interval restriction maps descend by
Proposition~\ref{prop:descent-temporal-restrictions}.

The trajectory sheaf comparisons are finite-limit isomorphisms. Their
pullbacks are the local trajectory sheaf comparisons, so they are
isomorphisms globally.
\end{proof}

\subsection{Descent of Segal-enhanced execution}
\label{subsec:descent-segal-enhanced-execution}

Let
\[
    \varphi:[m]\to[n]
\]
be any order-preserving map. It induces
\[
    a_\varphi^{\mathbb A}:
    A_n^\partial\to A_m^\partial.
\]
The corresponding simplicial structure map is regarded as a morphism over
\(A_n^\partial\):
\[
    N_+\operatorname{Exec}(\Phi)_n
    \to
    (a_\varphi^{\mathbb A})^\ast
    N_+\operatorname{Exec}(\Phi)_m.
\]

\begin{theorem}[Descent of Segal-enhanced execution]
\label{thm:descent-segal-enhanced-execution}
Under a nervewise cover, the augmented simplicial object
\[
    N_+\operatorname{Exec}(\Phi)
\]
satisfies effective interface descent degreewise. Its face and degeneracy
maps are the unique descents of the corresponding local maps.
\end{theorem}

\begin{proof}
Every nerve degree is an \(n\)-step execution object and descends along
\(p_n^\partial\).

For each
\[
    \varphi:[m]\to[n],
\]
the local simplicial structure map is a morphism in
\(\mathcal E/U_n^\partial\) from the degree-\(n\) object to the pullback of
the degree-\(m\) object along \(a_\varphi^{\mathbb U}\). The
Beck--Chevalley comparison associated with
\[
\begin{tikzcd}[column sep=large,row sep=large]
    U_n^\partial
        \arrow[r, "p_n^\partial"]
        \arrow[d, "a_\varphi^{\mathbb U}"']
    &
    A_n^\partial
        \arrow[d, "a_\varphi^{\mathbb A}"]
    \\
    U_m^\partial
        \arrow[r, "p_m^\partial"']
    &
    A_m^\partial
\end{tikzcd}
\]
identifies its target with the pullback of the corresponding global target.
It therefore descends uniquely.

The simplicial identities hold locally and are reflected globally by the
faithfulness of pullback along \(p_n^\partial\). The Segal maps are
finite-limit comparisons and are isomorphisms because their pullbacks are
isomorphisms.
\end{proof}

\begin{corollary}[Effective descent of derived temporal semantics]
\label{cor:stack-temporal-execution-semantics}
For every nervewise cover, the trajectory sheaves and augmented execution
nerves associated with a strict-system descent structure form effective
descent objects. Their reconstructions are respectively
\[
    \mathscr T_\Phi
\]
and
\[
    N_+\operatorname{Exec}(\Phi),
\]
where \(\Phi\) is the descended global system.
\end{corollary}

This statement concerns temporal semantics derived from strict systems. It
does not assert that every arbitrary labelled trajectory sheaf is the
trajectory semantics of a strict Mealy system.

\section{Temporal--interface interchange}
\label{sec:temporal-interface-interchange-rigorous}

\subsection{Temporal gluing over a common input-path base}
\label{subsec:temporal-gluing-recalled}

For every system \(\Phi\) and every \(m,n\in\mathbb N\), temporal gluing gives
\[
    \mathscr T_\Phi(m+n)
    \cong
    \mathscr T_\Phi(m)
    \times_{\mathscr T_\Phi(0)}
    \mathscr T_\Phi(n).
\]

To compare this pullback with interface descent, all terms must first be
viewed over the common input-path object
\[
    A_{m+n}^\partial.
\]

Let
\[
    \rho_m:
    A_{m+n}^\partial\to A_m^\partial
\]
be restriction to the initial \(m\)-segment, and let
\[
    \rho_n:
    A_{m+n}^\partial\to A_n^\partial
\]
be restriction to the terminal \(n\)-segment.

Let
\[
    \lambda_m:
    A_{m+n}^\partial\to A_0
\]
select the common boundary after \(m\) steps. The terminal-state map of the
initial segment and the initial-state map of the terminal segment then give
maps
\[
    \rho_m^\ast\mathscr T_\Phi(m)
    \to
    \lambda_m^\ast P,
\]
and
\[
    \rho_n^\ast\mathscr T_\Phi(n)
    \to
    \lambda_m^\ast P.
\]

Thus temporal gluing is the canonical isomorphism in
\(\mathcal E/A_{m+n}^\partial\):
\[
\begin{tikzcd}[column sep=large]
    \mathscr T_\Phi(m+n)
        \arrow[r, "\gamma_{m,n}^\Phi", "\cong"']
    &
    \rho_m^\ast\mathscr T_\Phi(m)
    \times_{\lambda_m^\ast P}
    \rho_n^\ast\mathscr T_\Phi(n).
\end{tikzcd}
\]

\subsection{Finite limits and interface reconstruction}
\label{subsec:finite-limits-descent-reconstruction}

For a regular epimorphism
\[
    e:E\twoheadrightarrow X,
\]
the pullback functor
\[
    e^\ast:\mathcal E/X\to\mathcal E/E
\]
is monadic and creates finite limits. Hence the equivalence between objects
over \(X\) and effective descent structures preserves finite limits.

\begin{proposition}[Finite-limit preservation of interface reconstruction]
\label{prop:finite-limit-preservation-interface-reconstruction}
Fix a nerve degree \(r\). Reconstruction along
\[
    p_r^\partial:
    U_r^\partial\twoheadrightarrow A_r^\partial
\]
preserves finite-limit diagrams in the slice
\[
    \mathcal E/A_r^\partial.
\]

Consequently, after transporting all terms to a common nerve degree,
reconstruction preserves:
\begin{enumerate}
    \item state-boundary pullbacks;
    \item path-object pullbacks;
    \item Segal comparisons;
    \item temporal gluing pullbacks.
\end{enumerate}
\end{proposition}

\begin{proof}
At the fixed degree \(r\), reconstruction is quasi-inverse to the effective
descent comparison induced by
\[
    (p_r^\partial)^\ast.
\]
Since this pullback functor creates finite limits, the descent equivalence
preserves them.
\end{proof}

\subsection{The interchange comparison}
\label{subsec:temporal-interface-interchange-comparison}

Let
\[
    p:\mathbb U\to\mathbb A
\]
be a nervewise cover, and let
\[
    \Phi_{\mathbb U}
\]
be a local system with a system descent structure. Let \(\Phi\) be its
descended global system.

There are two reconstruction procedures:
\begin{enumerate}
    \item glue adjacent temporal segments locally and then descend the
    resulting local execution;
    \item descend the temporal segments and their common boundary states, then
    glue globally.
\end{enumerate}

\begin{theorem}[Temporal--interface interchange]
\label{thm:temporal-interface-interchange-rigorous}
For every nervewise cover and every \(m,n\in\mathbb N\), the canonical
comparison between:
\begin{enumerate}
    \item temporal gluing followed by interface descent;
    \item interface descent followed by temporal gluing;
\end{enumerate}
is an isomorphism.

Equivalently, descent carries the local temporal gluing isomorphism to
\[
    \mathscr T_\Phi(m+n)
    \cong
    \mathscr T_\Phi(m)
    \times_{\mathscr T_\Phi(0)}
    \mathscr T_\Phi(n).
\]
\end{theorem}

\begin{proof}
View the temporal gluing diagram in the common slice
\[
    \mathcal E/A_{m+n}^\partial
\]
as
\[
    \mathscr T_\Phi(m+n)
    \xrightarrow{\cong}
    \rho_m^\ast\mathscr T_\Phi(m)
    \times_{\lambda_m^\ast P}
    \rho_n^\ast\mathscr T_\Phi(n).
\]

Pullback along
\[
    p_{m+n}^\partial:
    U_{m+n}^\partial\twoheadrightarrow A_{m+n}^\partial
\]
carries this diagram to the corresponding local temporal-gluing diagram.
This follows from trajectory base change and the Beck--Chevalley
isomorphisms for the prefix, suffix, and boundary maps.

By
Proposition~\ref{prop:finite-limit-preservation-interface-reconstruction},
the descent of the local pullback is the displayed global pullback. The
descended local gluing map and the canonical global gluing map have the same
pullback along \(p_{m+n}^\partial\). Faithfulness of pullback identifies them.
\end{proof}

\begin{remark}[Systems-theoretic interpretation of interchange]
\label{rem:systems-interpretation-interchange}
Suppose a long execution is divided into temporal segments and the input
environment is divided into overlapping interface patches.

One may:
\begin{enumerate}
    \item assemble the time segments inside each patch and then reconcile the
    resulting local runs across interfaces; or
    \item reconcile each time segment across interfaces and then concatenate
    the resulting global segments.
\end{enumerate}

The interchange theorem says that these procedures construct the same global
execution. Thus synchronization across interface descriptions is independent
of whether it is performed before or after temporal concatenation.
\end{remark}

\subsection{Segal interchange}
\label{subsec:segal-interface-interchange}

\begin{theorem}[Segal--interface interchange]
\label{thm:segal-interface-interchange}
Interface descent commutes with:
\begin{enumerate}
    \item insertion of identity executions;
    \item contraction of adjacent execution steps;
    \item all iterated face and degeneracy operations of the execution nerve.
\end{enumerate}
\end{theorem}

\begin{proof}
Each face or degeneracy operation is a morphism between execution objects
after its target has been pulled back to the source nerve degree. These maps
descend uniquely by
Theorem~\ref{thm:descent-segal-enhanced-execution}.
Their unit, associativity, and simplicial coherence equations hold locally
and are reflected globally.
\end{proof}

\subsection{Two-dimensional locality}
\label{subsec:two-dimensional-locality}

\begin{theorem}[Two-dimensional locality of execution]
\label{thm:two-dimensional-locality-execution}
The execution semantics of a strict Mealy system has two compatible
local-to-global structures:
\begin{enumerate}
    \item it is a sheaf in the discrete temporal interval variable;
    \item it satisfies effective reconstruction in the input-interface
    variable for nervewise covers.
\end{enumerate}
These reconstruction procedures commute.
\end{theorem}

\begin{proof}
Temporal sheafhood is the temporal gluing theorem of
Chapter~\ref{ch:temporal-trajectories-segal-execution}.
Effective interface reconstruction is
Theorem~\ref{thm:effective-interface-descent-temporal-trajectories}.
Their compatibility is
Theorem~\ref{thm:temporal-interface-interchange-rigorous}.
\end{proof}

\section{Formal examples and obstruction patterns}
\label{sec:formal-examples-interface-descent}

\subsection{Disconnected input interfaces}
\label{subsec:example-disconnected-input-interfaces}

Let
\[
    \mathbb A
    =
    \mathbb A_1\amalg\mathbb A_2.
\]
Then
\[
    \mathsf{Sys}_{\mathbb B}(\mathbb A)
    \simeq
    \mathsf{Sys}_{\mathbb B}(\mathbb A_1)
    \times
    \mathsf{Sys}_{\mathbb B}(\mathbb A_2),
\]
\[
    \mathsf{Ctr}_{\mathbb B}(\mathbb A)
    \cong
    \mathsf{Ctr}_{\mathbb B}(\mathbb A_1)
    \times
    \mathsf{Ctr}_{\mathbb B}(\mathbb A_2),
\]
and
\[
    \mathscr T_{\Phi_1\amalg\Phi_2}
    \cong
    \mathscr T_{\Phi_1}
    \amalg
    \mathscr T_{\Phi_2}.
\]

No overlap compatibility is required because there are no arrows connecting
the two input components.

\subsection{Object-refinement covers in \(\mathbf{Set}\)}
\label{subsec:example-object-refinement-cover}

Let \(\mathbb A\) be a small category in \(\mathbf{Set}\), and let
\[
    r:U_0\twoheadrightarrow A_0
\]
be a surjection.

Define a category \(\mathbb U_r\) by
\[
    (\mathbb U_r)_0:=U_0,
\]
and
\[
    (\mathbb U_r)_1
    :=
    U_0
    \times_{r,A_0,s_A}
    A_1
    \times_{t_A,A_0,r}
    U_0.
\]

An arrow is a triple
\[
    (u,a,v)
\]
such that
\[
    r(u)=s_A(a),
    \qquad
    r(v)=t_A(a).
\]
Projection defines
\[
    p_r:\mathbb U_r\to\mathbb A
\]
by
\[
    (p_r)_0=r,
    \qquad
    (p_r)_1(u,a,v)=a.
\]

\begin{proposition}[Object-refinement covers are nervewise]
\label{prop:object-refinement-covers-nervewise}
If \(r\) is surjective, then
\[
    p_r:\mathbb U_r\to\mathbb A
\]
is a nervewise cover.
\end{proposition}

\begin{proof}
An \(n\)-simplex of \(N\mathbb A\) is a composable string
\[
    a_1,\ldots,a_n
\]
with \(n+1\) vertices. Choose a lift under \(r\) of each vertex. These choices
produce a composable string in \(\mathbb U_r\) mapping to the original
string. Hence every
\[
    (p_r)_n
\]
is surjective.
\end{proof}

Thus strict systems and all finite trajectories descend along object
refinement.

However, \(p_r\) need not be behaviour-exact. A global arrow entering one
global object may have several lifts with different local source
representatives. The local residual slice can therefore contain distinctions
not present in the global slice, or fail to identify precisely the same
future experiments. Nervewise trajectory descent and exact residual semantics
are genuinely different requirements.

\subsection{A residual cover from a discrete fibration}
\label{subsec:example-residual-cover-discrete-fibration}

Let \(M\) be a monoid, let
\[
    \mathbb A=\mathbf BM,
\]
and let \(X\) be a nonempty right \(M\)-set.

Define a category
\[
    \mathbb U_X
\]
whose objects are elements of \(X\), and whose arrows are pairs
\[
    (m,x):
    x\cdot m\to x.
\]
Composition is
\[
    (m,x)\circ(n,x\cdot m)
    =
    (mn,x).
\]
Indeed,
\[
    (x\cdot m)\cdot n
    =
    x\cdot(mn),
\]
so the composite has source \(x\cdot(mn)\), target \(x\), and label \(mn\).

Projection defines
\[
    p_X:\mathbb U_X\to\mathbf BM
\]
by sending every object to the unique object and every arrow
\((m,x)\) to \(m\).

For every \(x\in X\) and every monoid arrow \(m\), there is a unique lifted
arrow with target \(x\). Hence \(p_X\) is a discrete fibration.

\begin{proposition}[Discrete-fibration residual cover]
\label{prop:discrete-fibration-residual-cover}
The functor
\[
    p_X:\mathbb U_X\to\mathbf BM
\]
is nervewise covering, universally behaviour-exact for every output category
\(\mathbb B\), and universally derivative-covering. Hence it is a nervewise
residually effective cover.
\end{proposition}

\begin{proof}
Nonemptiness of \(X\) implies surjectivity on objects. Every finite monoid
word has a unique lift ending at any chosen \(x\in X\), so every nerve map is
surjective.

More generally, let
\[
    q:\mathbb V\to\mathbb C
\]
be any discrete fibration. For every \(v\in V_0\), unique incoming-arrow
lifting induces an isomorphism
\[
    q_{/v}:
    \mathbb V/v
    \xrightarrow{\cong}
    \mathbb C/q_0v.
\]
Every arrow entering \(q_0v\) has a unique lift entering \(v\), and the same
uniqueness determines all morphisms in the slice. Residual precomposition is
therefore an isomorphism for every output category.

Likewise, the incoming-arrow comparison
\[
    \lambda_q:
    V_1
    \to
    V_0\times_{q_0,C_0,t_C}C_1
\]
is an isomorphism.

Pullbacks of discrete fibrations are discrete fibrations. Thus these slice and
incoming-arrow conclusions hold after every base change of \(p_X\).
Consequently, \(p_X\) is universally behaviour-exact and universally
derivative-covering.
\end{proof}

This example satisfies the system, trajectory, and residual reconstruction
conditions simultaneously.

\subsection{Pairwise compatibility without triple coherence}
\label{subsec:example-pairwise-no-triple-coherence}

Let
\[
    \{\mathbb U_i\to\mathbb A\}_{i=1}^3
\]
be a covering family. Suppose local systems
\[
    \Phi_i
\]
are equipped with pairwise overlap isomorphisms
\[
    \vartheta_{12},
    \qquad
    \vartheta_{23},
    \qquad
    \vartheta_{13}.
\]

If
\[
    \pi_{23}^\ast\vartheta_{23}
    \circ
    \pi_{12}^\ast\vartheta_{12}
    \neq
    \pi_{13}^\ast\vartheta_{13}
\]
on the triple overlap, then the local systems do not define a descent
structure.

Each pair of local models can then be reconciled, but no globally consistent
identification of their state spaces exists.

\subsection{Contract descent without implementation descent}
\label{subsec:example-contract-without-implementation-descent}

Suppose local residual contracts
\[
    V_i
\]
agree on overlaps. Contract descent determines a unique global contract
\[
    V.
\]

Choose local realizing systems
\[
    \Phi_i\models V_i.
\]
The systems \(\Phi_i\) need not admit overlap isomorphisms, even though the
contracts do. Therefore
\[
    \text{specification descent}
    \not\Rightarrow
    \text{descent of chosen implementations}.
\]

Conversely, a global system may exist while independently chosen local
contracts fail to agree on overlaps. This is a specification inconsistency
rather than an implementation inconsistency.

\subsection{Systems interpretations}
\label{subsec:systems-interpretations-interface-descent}

An input-interface cover may model:
\begin{enumerate}
    \item overlapping sensor views;
    \item distributed command or port sets;
    \item local operating modes;
    \item subsystem boundaries;
    \item local disturbance models;
    \item different protocol presentations;
    \item local controller information;
    \item overlapping physical and software models.
\end{enumerate}

For example, consider a controller with startup, normal, degraded, and
emergency modes. Different engineering teams may model overlapping
collections of modes.

In that interpretation:
\begin{enumerate}
    \item degree-zero coverage means every mode is represented;
    \item degree-one coverage means every command or mode transition is
    represented;
    \item degree-two coverage means every legally successive command pair is
    jointly represented;
    \item nervewise coverage means every finite command history is represented;
    \item behaviour exactness means each local model has exactly the correct
    future response profiles;
    \item derivative covering means every globally possible disturbance at a
    represented mode is included in the local model.
\end{enumerate}

Failure of descent may indicate:
\begin{enumerate}
    \item a missing global state variable;
    \item incompatible local control laws;
    \item a hidden interaction invisible on individual patches;
    \item pairwise agreement without triple coherence;
    \item a locally consistent but globally unrealizable implementation.
\end{enumerate}

\section{Comparison, limitations, and later directions}
\label{sec:comparison-limitations-interface-descent}

\subsection{Temporal and interface locality}
\label{subsec:comparison-temporal-interface-locality}

The two locality principles act on different variables.

Temporal locality fixes the input category and varies the execution interval:
\[
    [k]\to[n].
\]

Interface locality fixes the temporal shape and varies the input category:
\[
    F:\mathbb A'\to\mathbb A.
\]

The trajectory base-change theorem identifies their interaction, and the
interchange theorem states that the two local-to-global procedures commute.

\subsection{Hierarchy of exactness conditions}
\label{subsec:hierarchy-exactness-conditions}

The cover conditions have the following systems meanings:
\[
\begin{array}{c|c}
    \text{condition}
    &
    \text{structure visible locally}
    \\
    \hline
    p_0\text{ regular epic}
    &
    \text{all input contexts, state carriers, and unit actions}
    \\
    p_1\text{ regular epic}
    &
    \text{all one-step inputs, transitions, and outputs}
    \\
    p_2\text{ regular epic}
    &
    \text{all two-step interactions needed for the Mealy laws}
    \\
    p_n\text{ regular epic for all }n
    &
    \text{all finite input histories}
    \\
    p^\sharp\text{ invertible}
    &
    \text{exact preservation of residual future behaviour}
    \\
    \lambda_p\text{ regular epic}
    &
    \text{all global inputs applicable at represented contexts}
\end{array}
\]

Accordingly:
\begin{enumerate}
    \item system descent does not automatically imply residual descent;
    \item residual exactness does not automatically imply nervewise trajectory
    descent;
    \item behaviour exactness alone does not ensure complete derivative
    generation or safety checks;
    \item derivative covering alone does not reconstruct canonical residual
    behaviours.
\end{enumerate}

\subsection{Limitations}
\label{subsec:limitations-interface-descent}

The theory of this chapter is restricted to:
\begin{enumerate}
    \item strict Mealy systems and strict semantic homomorphisms;
    \item variation of the input interface while the output category remains
    fixed;
    \item covers internal to a Grothendieck topos;
    \item finite linear execution trajectories;
    \item residual contracts expressed as derivative-invariant subobjects.
\end{enumerate}

It does not yet provide:
\begin{enumerate}
    \item effective descent for general lax morphisms or simulations;
    \item simultaneous descent in both input and output interfaces;
    \item descent for hybrid or continuous-time trajectory sheaves;
    \item cubical, pomset, or concurrent execution descent;
    \item descent of feedback or traced closure;
    \item a classification of every behaviour-exact input functor;
    \item obstruction invariants for failed implementation descent;
    \item an independent characterization of all trajectory sheaves arising
    from strict Mealy systems.
\end{enumerate}

\subsection{Later directions}
\label{subsec:later-directions-interface-descent}

The relative program semantics of later chapters may be organized over varying
input interfaces. Its diamonds, boxes, and predicate transformers should
satisfy Beck--Chevalley laws over residually effective covers.

Polynomial response semantics may likewise be organized as an indexed family
of coalgebraic constructions. Behaviour exactness should then compare local
and global response profiles.

Concurrency may replace
\(\mathbf{Int}_{\mathrm{disc}}\) by cubical or partially ordered execution
shapes. The resulting theory should involve reconstruction simultaneously in
an interface variable and in a nonlinear execution-shape variable.

Feedback introduces another exactness problem: whether traced closure commutes
with input restriction and effective descent.

\section{Chapter synthesis}
\label{sec:interface-descent-chapter-synthesis}

\subsection{System descent}
\label{subsec:system-descent-summary}

\begin{theorem}[System descent theorem]
\label{thm:system-descent-summary}
For fixed output category \(\mathbb B\):
\begin{enumerate}
    \item strict Mealy systems form a pseudofunctor
    \[
        \mathsf{Sys}_{\mathbb B}:
        \mathsf{Inp}^{\mathrm{op}}
        \to
        \mathbf{Cat};
    \]
    \item extensive coproduct decompositions satisfy effective descent;
    \item every \(2\)-nerve cover
    \[
        p:\mathbb U\to\mathbb A
    \]
    induces an equivalence
    \[
        \mathsf{Sys}_{\mathbb B}(\mathbb A)
        \simeq
        \mathsf{Desc}_p
        \left(
            \mathsf{Sys}_{\mathbb B}
        \right);
    \]
    \item consequently,
    \[
        \mathsf{Sys}_{\mathbb B}
    \]
    is a stack for
    \[
        \tau_{\mathrm{sys}}.
    \]
\end{enumerate}
\end{theorem}

\begin{proof}
These are
Definition~\ref{def:input-indexed-system-pseudofunctor},
Theorem~\ref{thm:extensive-descent-strict-systems},
Theorem~\ref{thm:effective-descent-strict-mealy-systems}, and
Corollary~\ref{cor:stack-strict-mealy-systems}.
\end{proof}

\subsection{Residual descent}
\label{subsec:residual-descent-summary}

\begin{theorem}[Residual descent theorem]
\label{thm:residual-descent-summary}
For every input functor
\[
    F:\mathbb A'\to\mathbb A:
\]
\begin{enumerate}
    \item residual precomposition defines a strict semantic homomorphism
    \[
        F^\sharp:
        F^\ast
        \mathbf{Beh}^{\mathrm{can}}_{\mathbb A,\mathbb B}
        \to
        \mathbf{Beh}^{\mathrm{can}}_{\mathbb A',\mathbb B};
    \]
    \item semantic direct reindexing satisfies
    \[
        \operatorname{Beh}_{\mathbb A'}(F^\ast\Phi)
        =
        F^\bullet
        \operatorname{Beh}_{\mathbb A}(\Phi).
    \]
\end{enumerate}

For residually effective input functors:
\begin{enumerate}
    \setcounter{enumi}{2}
    \item residual contracts satisfy effective descent;
    \item generated derivative closure and safety commute with restriction;
    \item the derivative closure adjoint triple is indexed;
    \item the abstraction--realization adjunction is indexed.
\end{enumerate}
\end{theorem}

\begin{proof}
Strict residual restriction is
Theorem~\ref{thm:strictness-residual-restriction}.
The general abstraction formula is
Theorem~\ref{thm:behavioural-abstraction-arbitrary-restriction}.
Contract descent is
Theorem~\ref{thm:effective-descent-residual-contracts}.
Compatibility of generated closure and safety is
Theorems~\ref{thm:generated-closure-commutes-restriction} and
\ref{thm:safety-core-commutes-restriction}.
The indexed adjunctions are
Corollary~\ref{cor:indexed-derivative-adjoint-triple} and
Theorem~\ref{thm:indexed-abstraction-realization-adjunction}.
\end{proof}

\subsection{Local minimization}
\label{subsec:local-minimization-summary}

\begin{theorem}[Local minimization theorem]
\label{thm:local-minimization-summary}
For every residual interface cover:
\begin{enumerate}
    \item contract satisfaction is local;
    \item behavioural separatedness is local;
    \item semantic minimization commutes with restriction;
    \item generated residual theories commute with restriction;
    \item residual reachability is local;
    \item minimal residually generated separated realizations may be verified
    and reconstructed from compatible local restrictions.
\end{enumerate}
\end{theorem}

\begin{proof}
These are
Theorem~\ref{thm:locality-contract-satisfaction-rigorous},
Theorem~\ref{thm:separatedness-interface-local},
Theorem~\ref{thm:semantic-minimization-local-restriction},
Proposition~\ref{prop:generated-theories-local-restriction},
Theorem~\ref{thm:reachability-interface-local}, and
Corollary~\ref{cor:local-myhill-nerode-principle}.
\end{proof}

\subsection{Two-dimensional locality}
\label{subsec:two-dimensional-locality-summary}

\begin{theorem}[Temporal--interface locality theorem]
\label{thm:temporal-interface-locality-summary}
For strict Mealy systems:
\begin{enumerate}
    \item execution categories commute with input base change:
    \[
        \operatorname{Exec}(F^\ast\Phi)
        \cong
        \mathbb A'^{\mathrm{op}}
        \times_{\mathbb A^{\mathrm{op}}}
        \operatorname{Exec}(\Phi);
    \]
    \item temporal trajectory formation commutes with input restriction:
    \[
        \mathscr T_{F^\ast\Phi}
        \cong
        F^\sharp_{\mathrm{traj}}\mathscr T_\Phi;
    \]
    \item finite trajectory objects satisfy effective descent along nervewise
    covers;
    \item the full Segal-enhanced execution nerve satisfies effective
    descent;
    \item temporal gluing and interface descent commute.
\end{enumerate}
\end{theorem}

\begin{proof}
Execution base change is
Theorem~\ref{thm:execution-category-base-change}.
Trajectory base change is
Theorem~\ref{thm:trajectory-base-change}.
Trajectory and Segal descent are
Theorems~\ref{thm:effective-interface-descent-temporal-trajectories} and
\ref{thm:descent-segal-enhanced-execution}.
Interchange is
Theorems~\ref{thm:temporal-interface-interchange-rigorous} and
\ref{thm:segal-interface-interchange}.
\end{proof}

\begin{remark}[Final conceptual summary]
\label{rem:interface-descent-final-conceptual-summary}
Categorical automata are local in two independent variables.

They are local in time because finite trajectories form a sheaf on the
discrete path site.

They are local in the input interface because compatible local state machines
reconstruct a global state machine whenever the interface cover exposes all
input contexts, one-step actions, and two-step interactions required by the
Mealy laws.

The hierarchy of cover conditions has a direct operational meaning:
\[
    \text{degrees }0,1,2
    \quad\text{control reconstruction of the machine rules},
\]
\[
    \text{all nerve degrees}
    \quad\text{control direct reconstruction of finite runs},
\]
\[
    \text{behaviour exactness}
    \quad\text{ensures that local futures are exactly global futures
    restricted},
\]
and
\[
    \text{derivative covering}
    \quad\text{ensures that no globally possible input is hidden from a local
    closure or safety calculation}.
\]

Under these hypotheses, local verification is complete, local minimization
agrees with restricted global minimization, and temporal assembly commutes
with interface assembly.

Hence the same global execution may be reconstructed either by first
assembling its temporal segments and then integrating its interface-local
descriptions, or by performing those two constructions in the opposite order.
\end{remark}